\artikelfoto{Van basisonderwijs naar secundair onderwijs}{}{Johan Deprez \\ Michèle Dexters \\ Els Van Emelen}{img/coverUW4002.jpg}

% TODO: beter met fototje

\startcontents[inhoudloep]
\begin{multicols}{2}

\inhoudstafelloep

\section{Inleiding}\label{inleiding}

Voor leerlingen gebeurt het leren van wiskunde in een heel lange leerlijn, startend vanaf de kleuterschool (en thuis) over de basisschool naar het secundair onderwijs en vaak ook nog verder naar het hoger onderwijs.
Doorheen die lange leerlijn krijgen ze lesgevers die per onderwijsniveau op een andere manier opgeleid zijn, elk met hun eigen cultuur en vaak ook met onderling weinig contact.
De overgang van een onderwijsniveau naar het volgende verloopt niet altijd even vlot.
In deze loep focussen we op de overgang tussen basisonderwijs en secundair onderwijs.
Onderwijzers weten niet altijd hoe met wiskunde verder gegaan wordt in het secundair onderwijs en wiskundeleraren uit de eerste graad van het secundair onderwijs zijn niet altijd op de hoogte van hoe wiskundige inhouden in het basisonderwijs behandeld werden.
Op die manier ontstaan er breuken in de leerlijn, wat niet goed is voor de leerlingen.

We werken in deze loep exemplarisch door de overgang in detail te bekijken voor een beperkt aantal onderwerpen.
Daarbij onderscheiden we twee verschillende situaties.
Als je eindtermen en leerplannen van het basisonderwijs en de eerste graad van het secundair onderwijs bekijkt, dan zie je veel onderwerpen terugkomen.
Dat is bijvoorbeeld het geval met meetkundige figuren, verhoudingen en bewerkingen.
Je zou er voor deze onderwerpen van uit kunnen gaan dat het puur over herhaling of opfrissing gaat.
Dan laat je echter kansen liggen.
Je kunt dezelfde onderwerpen in basisonderwijs en secundair onderwijs immers vanuit een verschillend perspectief bekijken, waarbij je bewust een trap naar (iets) meer abstractie inbouwt.
Wat we uitwerken in \autoref{meetkunde} over meetkunde is bijvoorbeeld van deze aard.

Bij andere onderwerpen zie je juist wel een duidelijk verschil tussen basisonderwijs en secundair onderwijs.
Denk hierbij bijvoorbeeld aan de start van de algebra.
Toch zagen we bij het voorbereiden van de loep dat er ook hier aanknopingspunten zijn met wat in de basisschool gebeurt.
Zo lossen leerlingen in de basisschool vraagstukken op rond `ongelijke verdeling' en wellicht zijn onze lezers vertrouwd met de `puntsommen' of `vleksommen' uit de basisschool.
Zo'n punt of vlek vormt in feite een informele voorloper voor de onbekende die we in het secundair onderwijs door een letter voorstellen.
Bij dit soort onderwerpen pleiten we ervoor om expliciet het verband te leggen tussen de nieuwe werkwijze (algebraïsch oplossen van een vergelijking) en wat de leerlingen vroeger geleerd hebben.
Zeg dus niet: vergeet wat je vroeger geleerd hebt, maar vertrek juist vanuit wat de leerlingen in de basisschool geleerd hebben en bouw daar op verder om dan de overgang te maken naar de nieuwe werkwijze.
Maak ook duidelijk waar de nieuwe manier van werken superieur is en durf ook erkennen wanneer ze eigenlijk niet nodig is.
We werken deze start van de algebra hier niet verder uit, want we hebben dit al eerder gedaan in de loep van UW29/1.
Nu werken we andere voorbeelden uit.
Percentrekenen en verhoudingstabellen (in \autoref{percent}) en coördinaten (in \autoref{getallen}) zijn daar voorbeelden van.

We hebben elke paragraaf volgens dezelfde structuur uitgewerkt.
We schetsen eerst hoe dit onderwerp in het basisonderwijs en in het secundair onderwijs aan bod komt.
Daarna geven we aan waar we een breuk of gemiste kansen zien en doen we voorstellen om de overgang te verbeteren.

Omdat het voor de loep van belang is, stellen we de auteurs even voor.
Johan is verantwoordelijk voor vakdidactiek wiskunde in de educatieve master van KU Leuven en bekijkt de overgang tussen beide niveaus dus vanuit het perspectief van het secundair onderwijs.
Michèle is gastauteur voor deze loep en is net als Els lector in de educatieve bachelor basisonderwijs van UCLL in Diepenbeek.
Michèle en Els voerden de afgelopen jaren onderzoek uit binnen het expertisecentrum Art of Teaching naar leerlijnen en curricula en de manier waarop die door leraren in de basisschool gekend zijn en gebruikt worden.
Op dit moment maakt Michèle ook deel uit van een lopend onderzoek met als uiteindelijk doel het ontwerpen van een online game dat leerlingen uit de eerste graad van het secundair onderwijs de mogelijkheid geeft om inhouden uit de basisschool te remediëren of verder in te oefenen, om zo de nodige basiskennis te verstevigen.
Overigens, scholen die willen meewerken aan dit onderzoek of die interesse betonen voor (het uittesten van) de game, mogen contact opnemen met Michèle (\href{mailto:michele.dexters@ucll.be}{\nolinkurl{michele.dexters@ucll.be}}).

We sluiten deze inleiding af met twee disclaimers.
Om zo'n breuk in de overgang te verhelpen, zijn er aan beide kanten van de breuk inspanningen nodig.
Omdat ons doelpubliek uit leraren secundair onderwijs bestaat, focussen we hier in Uitwiskeling op wat vanuit het secundair onderwijs kan gebeuren om hier beter mee om te gaan.
Natuurlijk hebben we ook suggesties voor de collega's uit de basisschool.
Ook daar liggen immers kansen om beter in te spelen op wat de toekomst voor hun leerlingen brengt.
Deze kant van het verhaal vind je niet in onze loep.
Daarvoor plannen we een gelijkaardig artikel voor een tijdschrift dat zich richt op leraren basisonderwijs.

Zoals steeds bij Uitwiskeling is deze loep geschreven door auteurs die dit grotendeels in hun vrije tijd doen.
We beperken ons tot een aantal onderwerpen waarover we, denken we, iets zinvols te vertellen hebben.
We hebben voor die onderwerpen een aantal handboeken geraadpleegd en we steunen op onze kennis van het werkveld, die uiteindelijk slechts fragmentarisch is.
De loep is dus niet het resultaat van een systematisch wetenschappelijk onderzoek over de overgang tussen basisonderwijs en secundair onderwijs (al vinden we dat dit best wel aan de orde is!).

\section{Percentrekenen en verhoudingstabellen}\label{percent}

\subsection{Leerstof in het basisonderwijs}

In de derde graad van de lagere school lossen leerlingen vraagstukken met procenten op m.b.v. een verhoudingstabel.
Dat zien we bijvoorbeeld in \autoref{fig:vaderparfum}.
\end{multicols}

\afbeelding[12cm]{img/loepMichele/figuurm1.png}{Een opgave over percentrekenen (Uit: \emph{Reken Maar 6}, blok 7, les 77)}{fig:vaderparfum}

\begin{multicols}{2}
De redenering gaat als volgt:

\begin{itemize}
\item Bij vraagstukken over korting spelen de begrippen (of grootheden) `oude prijs', `korting' en `nieuwe prijs' mee.
Deze begrippen behoren tot de kennis die leerlingen hebben over dit soort vraagstukken en worden in principe altijd op die manier benoemd, ook al worden er in het vraagstuk andere benamingen gebruikt.
In het bovenstaande vraagstuk wordt met de `normale prijs' uiteraard de `oude prijs' bedoeld.
\item Het kortingspercentage wordt altijd uitgedrukt ten opzichte van de oude prijs (= het geheel of de referentie).
\item 40 \% korting betekent dat wanneer de oude prijs \EUR{100} was, je \EUR{40} korting krijgt en de nieuwe prijs dus \EUR{60} is.
\item De nieuwe prijs is echter geen \EUR{60}, maar wel \EUR{36}.
Om van \EUR{60} naar \EUR{36} te gaan, kun je delen door 10 en vermenigvuldigen met 6.
(Idealiter staan er pijlen met deze bewerkingen boven de verhoudingstabel genoteerd.)
\item De oude prijs was dan ook geen \EUR{100}.
De nieuwe en de oude prijs zijn recht evenredige grootheden.
Om de oude prijs te berekenen, mag je dus dezelfde bewerking toepassen (delen door 10 en vermenigvuldigen met 6).
\[
\EUR{100} : 10 = \EUR{10} \text{ en } 6 \times \EUR{10} = \EUR{60}.
\]
Merk op dat we \(6 \times \EUR{10}\) schrijven en niet \(\EUR{10} \times 6\).
De vermenigvuldiging is wel commutatief, maar in de lagere school leer je in eerste instantie dat vermenigvuldigen overeenkomt met `groepjes nemen' of `keer'.
`6 keer 10 euro' heeft betekenis, '10 euro keer 6' niet.
\end{itemize}

\subsection{Leerstof in het secundair onderwijs}

In de eerste graad van het secundair onderwijs komen dezelfde soort opgaven aan bod, maar \autoref{fig:laptop} toont een meer abstracte manier van oplossen.

\end{multicols}
\afbeelding[14cm]{img/loepMichele/figuurm2.png}{Percentrekenen in het secundair (Uit: \emph{Nando 2}, H9)}{fig:laptop}
\begin{multicols}{2}

\subsection{Hoe bouw je de brug?} 
We zetten de redeneringen en berekeningen uit beide oplossingsmethoden
naast elkaar en zoeken naar overlap en aanknopingspunten (zie figuur \autoref{fig:tabelopl}).

Je merkt dat de stap van een verhoudingstabel naar een evenredigheid niet zo groot is.
Als je in het laatste vakje van de BTW nu nog een \(x\) plaatst, zie je in de verhoudingstabel als het ware de evenredigheid verschijnen.
Je slaat de brug tussen het basisonderwijs en het secundair onderwijs door deze overgang met leerlingen te bespreken.
Je kunt ook expliciet aangeven dat de berekeningen die je maakt wanneer je een evenredigheid via een vergelijking oplost, dezelfde zijn als de bewerkingen die er bij de pijlen in het verhoudingsschema verschijnen.
In de basisschool besteed je, zoals eerder vermeld, wel nog aandacht aan de functie van vermenigvuldiger en vermenigvuldigtal: je neemt 1452 keer 1 euro, vandaar de schrijfwijze `\(1452 \times \)' in de verhoudingstabel.
In het secundair onderwijs ga je ervan uit dat leerlingen hun berekeningen tijdelijk los kunnen zien van de context of m.a.w. abstractie van de context kunnen maken.
Daardoor is het dan wel mogelijk om eigenschappen, in dit geval commutativiteit, toe te passen.
Dan pas kun je `\(\times 1452\)' schrijven.

Na een periode van overlap, waarbij beide werkwijzen naast elkaar gebruikt worden, laat je de leerlingen overstappen op de meer abstracte notatie en denkwijze.
Die is dan echter gestoeld op een concreter, inzichtelijk begrip van de situatie, waar ze, in geval van nood, op kunnen terugvallen.

We kunnen nog een stap verder gaan.
Daarvoor hernemen we de verhoudingstabel en voegen we er een lijn aan toe:


\afbeelding{img/loepMichele/figuurm4.jpg}{Aangevulde verhoudingstabel}{fig: aangevuld}

In plaats van `horizontaal' te redeneren `binnen' elk van de grootheden, kunnen we ook `verticaal' redeneren.
We zoeken dan het verband tussen de grootheden.

Zo zien we dat
\begin{align*}
    \text{brutoprijs} &=\text{nettoprijs} + \num{0.21} \times \text{nettoprijs}\\
    &= \num{1.21} \times \text{nettoprijs}
\end{align*}

Op die manier zetten we de stap naar het beschrijven van een evenredigheid met behulp van een functie.
Ook het belangrijke begrip evenredigheidsfactor komt hier naar voor.
Hier ligt dus eveneens een kans om te vertrekken vanuit wat leerlingen uit het basisonderwijs meebrengen en behoedzaam abstractiestappen te zetten.
In het vervolg van het secundair onderwijs wordt hier dan op verder gebouwd.
De speciale eerstegraadsfuncties die we hier gebruiken leggen de basis voor de algemene eerstegraadsfuncties en de tweedegraadsfuncties in (meestal) de tweede graad en een heel arsenaal aan andere functies in de derde graad.


\end{multicols}
\afbeelding[14cm]{img/loepMichele/figuurm3.jpg}{Vergelijkende tabel met oplossingsmethoden uit de lagere en de secundaire school}{fig:tabelopl}

\begin{multicols}{2}

\subsection{Mind the gap}

Er zijn een aantal factoren die dit abstraheringsproces bemoeilijken.

\subsubsection*{Onvoldoende inzicht in de betekenis en de asymmetrie van \%}

De fout die bij percentrekenen het vaakst voorkomt, heeft met inzicht in de betekenis van percenten te maken.
Een percent staat nooit op zichzelf en wordt altijd genomen van een gedefinieerd of afgesproken geheel.
We noemen dit de referentie.

In het voorbeeld van de BTW zijn veel leerlingen (en volwassenen) geneigd om 21 \% te nemen van \EUR{1452} wanneer ze de BTW zoeken.
Deze leerlingen staan niet stil bij de vraag wat de referentie is.
Hoewel dit inzicht dus ook al onderwerp van gesprek is in de basisschool, kan het niet voldoende herhaald worden, telkens wanneer leerlingen met percenten in aanraking komen.
Zeg dus (in eerste instantie) niet zomaar dat de BTW 21 \% bedraagt, maar wel 21 \% \textbf{van de netto prijs}.
De ene lagere school is de andere niet en ook de gebruikte methodes en werkboeken verschillen.
Sommige leerlingen zullen meer geoefend zijn om zich af te vragen wat de referentie is dan andere.

\subsubsection*{Versnippering in de leerwerkboeken in het secundair onderwijs}

Leerwerkboeken uit de eerste graad van het secundair onderwijs maken een opdeling in hoofdstukken waardoor het verband tussen een aantal essentiële begrippen (verhouding, evenredigheid, percentage) te weinig zichtbaar wordt.



\end{multicols}
\afbeelding[15cm]{img/mindthegap.jpg}{}{}
\begin{multicols}{2}

In hetzelfde werkboek waar de oefening over BTW gegeven wordt, vinden we enkele hoofdstukken eerder (!) nog een andere opdracht over percentrekenen terug (zie \autoref{fig:versnippering}).

\end{multicols}
\afbeelding[14cm]{img/loepMichele/figuurm5.png}{Eerdere percentopgave in hetzelfde handboek (Uit: \emph{Nando 2}, H7)}{fig:versnippering}
\begin{multicols}{2}

Hoewel de titel van het hoofdstuk doet vermoeden dat er met percenten én verhoudingen gewerkt wordt, legt men in het theoretisch gedeelte enkel uit dat het nemen van een percentage overeenkomt met het vermenigvuldigen met een breuk met noemer 100:
\[
30 \% \text{ van } x = \frac{30}{100} \cdot x = \num{0.3} \cdot x
\]

De verhoudingstabel uit de basisschool is zeer moeilijk te vertalen naar de bovenstaande werkwijze.

Er is niks op tegen om op termijn naar deze verkorte berekening van percenten te evolueren.
Zoals we hierboven aangaven, is dit echter pas een laatste stap in de abstrahering.
Vertrekkende vanuit verhoudingstabellen die gekend zijn uit de lagere school, kan er eerst horizontaal geredeneerd worden om zo naar evenredigheden oplossen m.b.v. vergelijkingen te gaan.
Daarna kun je ook verticaal redeneren om zo te evolueren naar functies.

Bepaalde keuzes die in werkboeken gemaakt worden, zitten dit soort logische kennisopbouw en de mogelijkheid tot het leggen van de nodige verbanden in de weg.
Het loont dus zeker de moeite om de voorstellen van de methodes met een kritische blik te bekijken en eventueel eigen keuzes te maken.

\section{Het verband tussen breuken en kommagetallen}\label{breuk}

\subsection{Leerstof in het basisonderwijs}
Vanaf het vierde leerjaar leren kinderen om kommagetallen om te zetten naar breuken en omgekeerd.
In \autoref{fig:breukexample} vind je een typische opgave terug.

\end{multicols}

\afbeelding[13.5cm]{img/loepMichele/figuurm6.png}{Typische opgave omzetten breuken en kommagetallen (uit: \emph{Reken Maar 4}, blok 12, les 134)}{fig:breukexample}

\begin{multicols}{2}
    
Er zijn verschillende manieren om dit verband te vinden.

\subsubsection*{Het kommagetal lezen volgens zijn rangen}\label{het-kommagetal-lezen-volgens-zijn-rangen}

In opdracht a uit \autoref{fig:breukexample} zijn de getallen zo gekozen dat het volstaat om het kommagetal te lezen volgens zijn rangen.
Het getal \(\num{0.3}\) kun je bijvoorbeeld lezen als 0 eenheden (E) en 3 tienden (t) of `drie tienden', wat rechtstreeks leidt tot de breuk \(\frac{3}{10}\).

De eerste twee opgaven uit opdracht b kun je op analoge wijze vinden.
De breuk \(\frac{15}{100}\) lees je als `vijftien honderdste', wat je ook kunt noteren als 15h en zo vind je \(\num{0.15}\).

\subsubsection*{De breuk omzetten naar een decimale breuk}

Bij de laatste opgaven wordt dat echter moeilijker.
Je moet de breuk eerst omzetten naar een decimale breuk (een breuk met als noemer een macht van 10) om dezelfde strategie te kunnen blijven volgen.
\(\frac{4}{5} = \frac{8}{10}\) , die je kunt lezen als `acht tiende' of 8t.
\(\frac{4}{5}\) is dus gelijk aan 8t of \(\num{0.8}\).

Een breuk die je niet kan omzetten naar een decimale breuk, leren de meeste leerlingen in de lagere school nog niet omzetten naar een kommagetal.

\subsection{Leerstof in het secundair onderwijs}

In het eerste jaar van het secundair onderwijs wordt een breuk formeel gedefinieerd.
In \autoref{fig:breukformeel} zie je een fragment.

\end{multicols}
\afbeelding[14cm]{img/loepMichele/figuurm7beter.png}{Definitie van een breuk (uit: \emph{Delta Nova 1b}, p. 9)}{fig:breukformeel}
\begin{multicols}{2}

Op basis van deze definitie kun je, naast de twee soorten breuken die hierboven al besproken werden, nog een derde soort aan de oefeningen toevoegen.
Een breuk die je niet kunt herleiden tot een decimale breuk, kun je immers wel via een deling omzetten naar een kommagetal.
Zo'n kommagetal heeft dan een oneindig aantal decimalen met een repeterend deel, de periode genaamd.
\[
\frac{2}{3} = 2\ :3 = \num{0.66}\ldots
\]

\subsection{Hoe bouw je de brug?}

Dat je een breuk kunt opvatten als een deling wordt in de meeste lagere scholen niet echt benadrukt.
Dit inzicht heb je nodig om verder te kunnen met het concept `rationaal getal' en zo o.a. alle soorten breuken, ook de niet-decimale, om te zetten naar een kommagetal.

Via getallenassen is het mogelijk om inzichtelijk aan te brengen dat een breuk ook kan opgevat worden als een deling.

\subsubsection*{Punten op de getallenas benoemen als een breuk}
We vragen leerlingen eerst om bepaalde plaatsen op de getallenas te benoemen met een breuk op een getallenas (zie \autoref{breukaflezen}).
In dit voorbeeld moeten leerlingen inzien dat de eenheid, de afstand tussen 0 en 1, verdeeld is in 4 gelijke delen.
De eerste breuk die ze moeten noteren, staat 3 van die deeltjes van de 0 verwijderd en slechts eentje van de 1.
Hij duidt dus 3 van de 4 gelijke delen aan.
Het gaat om de breuk \(\frac{3}{4}\).
\end{multicols}
\afbeelding[14cm]{img/loepMichele/figuurm8.jpg}{Breuken aflezen op getallenas}{breukaflezen}

\begin{multicols}{2}


Aangezien we ons op de getallenassen die we al kennen met natuurlijke getallen, niet beperken tot het interval tussen 0 en 1, vragen we ook om breuken te noteren die groter zijn dan het geheel, de zogenaamde `onechte breuken'.
De tweede breuk die we moeten noteren staat 5 deeltjes van de 0 verwijderd of een deeltje verder dan de 1.
Het gaat om de breuk \(\frac{5}{4}\) of om het gemengd getal 1 en \(\frac{1}{4}\).

\subsubsection*{Breuken op een getallenas plaatsen}
De omgekeerde oefening, waarbij leerlingen gegeven breuken op een as moeten plaatsen, is een volgende stap. In \autoref{fig:breukenplaatsen} vind je een dergelijke opgave.
Idealiter laat je hen zelf een passende verdeling aanbrengen of m.a.w. de as ijken, zoals we verder bespreken in \autoref{getallen}.
Ze leren om het geheel zo te kiezen, dat je het gemakkelijk kunt verdelen volgens de waarden die de noemers aangeven.
Voor de breuken hieronder, die niet toevallig allemaal herleidbaar zijn tot een decimale breuk, is een verdeling van de eenheid in 10 gelijke delen aangewezen.

\end{multicols}
\afbeelding[14cm]{img/loepMichele/figuurm9.jpg}{Breuken plaatsen op getallenas}{fig:breukenplaatsen}
\begin{multicols}{2}

Om de breuk \(\frac{9}{5}\) op de juiste plaats te zetten, verdeel je de eenheid in 5 gelijke delen (twee streepjes per deel) en neem je 9 van zulke delen.

\subsubsection*{Invoeren van kommagetallen}
Nadat je deze opdracht hebt gemaakt, is het interessant om een nieuwe probleemstelling te geven: het is helemaal niet zo handig om breuken op een getallenas te schrijven.
Ze zien er zo anders uit dan de getallen die we tot dan toe gewoon zijn en het zou eenvoudiger of eenduidiger zijn om een schrijfwijze te hanteren die nauwer aansluit bij die van de natuurlijke getallen.
Bij natuurlijke getallen zijn we het gewoon om per 10 te denken.
We leerden ooit dat 10 eenheden konden worden ingewisseld voor een tiental of omgekeerd, dat een eenheid 1 van de 10 gelijke delen is van een tiental.
We kunnen dezelfde redenering doorvoeren wanneer we de eenheid verder verdelen.
Als we die in 10 gelijke delen verdelen, dan kunnen we zo'n deeltje `een tiende' noemen, naar analogie met de breuk die daarmee overeenkomt op de getallenas.
De nieuwe schrijfwijze kunnen we dan baseren op het plaatswaardesysteem dat in de lagere leerjaren naar links werd uitgebreid (E, T, H, D, enz.) en nu ook naar rechts (t, h, d).
Het enige probleem is dan dat je wel een manier moet vinden om aan te geven waar de eenheid zich bevindt.
Tot nu toe kon je ervan uitgaan dat het meest rechtse cijfer, het cijfer van de eenheden was.
We voeren vanaf nu het symbool van de komma in om aan te geven waar de scheiding tussen de eenheid en de decimalen precies zit (zie \autoref{fig:invoerenkomma}).

Omdat we op zoek zijn gegaan naar een eenvoudigere of meer eenduidige schrijfwijze voor breuken op een getallenas, hebben we de kommagetallen ingevoerd.
Hierdoor hebben ze automatisch een link met breuken.
Wanneer je dan bijvoorbeeld de breuk \(\frac{4}{5}\) moet omzetten naar een kommagetal, zoals in de b-opdracht van \autoref{fig:breukexample}, ga je niet op zoek naar een gelijkwaardige decimale breuk omdat je dan `hoort' om welk kommagetal het gaat (de meer intuïtieve benadering).
Je gaat op zoek naar de gelijkwaardige decimale breuk \(\frac{8}{10}\), omdat je begrijpt dat kommagetallen opgebouwd zijn uit `tiende delen' van de eenheid.
\(\frac{8}{10}\) staat voor 8 van de 10 gelijke delen van de eenheid of dus ook 8t of \(\num{0.8}\).

\end{multicols}
\afbeelding[14cm]{img/loepMichele/figuurm10BETER.png}{Het invoeren van de komma bij decimale getallen}{fig:invoerenkomma}
\begin{multicols}{2}

\subsubsection*{Van niet-decimale breuk naar een kommagetal of deling}

Van zodra leerlingen begrijpen wat het echte wiskundige verband is tussen breuken en kommagetallen, kan deze voorkennis ook gebruik worden om een volgende moeilijke stap te zetten.

Om de breuk \(\frac{1}{3}\) op een getallenas te zetten, kun je de eenheid verdelen in 3 gelijke delen.
Je zet de breuk dan bij het eerste streepje rechts van de 0.
Wanneer je deze breuk echter wil omzetten naar een kommagetal, botsen we op het probleem dat je niet verder kunt verdelen tot je een tiendelige verdeling hebt of m.a.w. je vindt geen gelijkwaardige breuk met als noemer een macht van 10.

Wanneer leerlingen het verband tussen breuken en kommagetallen echt begrepen hebben, kun je hier het nieuwe inzicht en een nieuwe graad van abstractie invoeren, nl. dat een breuk ook kan opgevat worden als een deling.

De breuk \(\frac{1}{3}\) betekent dat je de eenheid, zijnde 1E, verdeelt in 3 gelijke delen of m.a.w. deelt door 3.
Eerder leerden we al dat wanneer we `te kort hebben' om te delen, we kunnen inwisselen voor een lagere rang.
Wanneer je 500 wilt delen door 25 dan merk je dat je met 5H niet ver komt.
Wanneer je de 5H inwisselt voor 50T, dan vind je dat je 2T uitkomt bij deling door 25 of m.a.w. 500 : 25 = 2T = 20.
Datzelfde principe kun je nu ook gebruiken.
1E kunnen we niet delen door 3.
We wisselen 1E in voor 10t.
Wanneer je 10t verdeelt over 3, dan krijgt ieder in elk geval al 3t.
De rest van 1t kun je weer inwisselen voor 10h, waarvan je er ieder ook 3 kunt geven, enz.
Je vindt dat \(\frac{1}{3} = 1 : 3 = \num{0.33}\ldots\).

Voor de volledigheid kun je deze ontdekking nog controleren voor de eerder gekende omzettingen.
\[
\frac{3}{10} = 3 : 10 = 30 t : 10 = 3 t = \num{0.3}
\]
\[
\frac{4}{5} = 4 : 5 = 40 t : 5 = 8 t = \num{0.8} = \frac{8}{10}
\]
Leerlingen krijgen zo een allesomvattend beeld van getallen en de verbanden tussen de verschillende notaties.
Een breuk kun je opvatten als een deling, waardoor het verband tussen breuken en kommagetallen nog scherper wordt gemaakt.

\subsection{Mind the gap}

In de meeste curricula uit het lager onderwijs kiest men ervoor om breuken en kommagetallen apart van elkaar aan te brengen.
Vanuit de leefwereld komen kinderen met o.a. geld in aanraking, waardoor jonge kinderen ook op school al kommagetallen bestuderen, zij het op een intuïtieve manier.
Breuken brengt men systematischer aan vanaf het derde leerjaar, in eerste instantie om een deel van een continu geheel, zoals een taart, te benoemen.
In het vierde leerjaar moeten kinderen de overstap maken naar breuk als rationaal getal en dan gaat men ineens het verband leggen met de - reeds gekende - kommagetallen.
Een noodzaak of probleemstelling ontbreekt echter meestal en de aanpak blijft dan ook beperkt tot het lezen van de breuk en het kommagetal en zo dus het `horen' van het verband.
Een aanpak die jammer genoeg niet de nodige voorkennis opbouwt om het moeilijkere inzicht van de breuk als deling te kunnen vatten.

Hoewel twaalfjarigen dus al heel wat weten over breuken en kommagetallen, is de voorkennis waarop je kunt voortbouwen wellicht verschillend.
Wanneer je op een abstracter niveau met hen over getallen wilt spreken, is het zinvol om het bovenstaande denkproces vanaf het invoeren van de getallenassen (opnieuw) te doorlopen.
Op die manier sluit je met meer zekerheid aan bij hun voorkennis en kunnen ze de sprong naar de volgende graad van abstractie gefundeerd maken.

\section{Eigenschappen van bewerkingen in functie van het rekenen met letters}\label{bewerk} % michele 

\subsection{Leerstof in het basisonderwijs}
De basisschool is de plek waar leerlingen hun eerste ervaringen opdoen met bewerkingen.
Zo leren leerlingen onder andere op verschillende manieren aftrekken.
In de meeste scholen wordt er gekozen voor een standaardmethode, waarbij het aftrektal in z'n geheel wordt behouden en de aftrekker gesplitst wordt in rangen.
Op deze manier kun je voor elke opgave een antwoord vinden en voor de meeste opgaven is dit zelfs de meest efficiënte manier.
Bijvoorbeeld:
\begin{align*}
    57-28 &= 57-20-8\\
    &= 37-8\\
    &= 29
\end{align*}

In de tweede en de derde graad van de basisschool leren kinderen ook een aantal handigere methodes zoals:
\begin{itemize}
    \item aftrekkers van plaats wisselen:
\begin{align*}
    138 - 45 - 28 &= 138 - 28 - 45 \\
    &= 110 - 45 \\
    &= 70 - 5 = 65 
\end{align*}

    \item  compenseren:
\[
    1357 - 99 = 1357 - 100 +1
\]
    \item gebruik maken van de aftrekkingshalter:
\[
    1357 - 99 = (1357 + 1) - (99 + 1)
\]
\end{itemize}
Het begrip aftrekkingshalter wordt gebruikt om aan te geven dat het resultaat van een aftrekking niet verandert wanneer je bij zowel aftrektal als aftrekker hetzelfde getal optelt of aftrekt.
In sommige scholen kiest men eerder voor compenseren, waarbij je dat wat je te veel hebt afgetrokken, er logischerwijze weer moet bijtellen.

Ook voor de andere bewerkingen komen er standaardmethodes en een aantal handige methodes aan bod in de lagere school.

Van in het eerste leerjaar onderzoeken kinderen ook enkele `spelregels' bij deze bewerkingen.
Zo leren ze dat je, in tegenstelling tot bij de aftrekking en de deling, bij de optelling en de vermenigvuldiging de termen en factoren altijd `van plaats mag wisselen' en dat je bij deze bewerkingen ook altijd mag `schakelen'. 
De namen commutativiteit en associativiteit worden niet gebruikt.

In \autoref{fig:commu} vind je een fragment uit een les in het vijfde leerjaar.
De leerlingen oefenen op twee handige rekenmethodes om te vermenigvuldigen.
Je ziet hoe deze `regeltjes', of beter gezegd eigenschappen, in het werkboek worden weergegeven.
Merk op: distributiviteit noemt men `splitsen en verdelen' in de ene richting en `factoren samennemen' in de andere richting.

\end{multicols}
\afbeelding[14cm]{img/loepMichele/figuurm11.png}{Rekenregels in het basisonderwijs (uit: \emph{Wiskanjers 5}, blok 3, les 4)}{fig:commu}
\begin{multicols}{2}

Ten slotte leren kinderen ook de volgorde van bewerkingen, inclusief het gebruik van haakjes, maar uiteraard zonder machten en wortels.

\subsection{Leerstof in het secundair onderwijs}
In het eerste leerjaar van het secundair onderwijs gaat het snel.
Na een herhaling van de vier hoofdbewerkingen en de verbanden daartussen (optellen en aftrekken als omgekeerde bewerkingen, som van gelijke termen schrijven als een product, enz.), staan al vlug machten en wortels als volgende stap op het programma.
De eerder gekende volgorde van bewerkingen wordt uitgebreid met deze nieuwe `bewerkingen'.

Verder wordt er in het secundair onderwijs systematisch gerekend met breuken en met negatieve getallen, daar waar die maar in (heel) beperkte mate aan bod kwamen in de basisschool.
Dat is een moment waarop veel leerlingen de draad kwijt geraken, onder andere door de vele regeltjes die hiervoor worden afgeleid.

Op het einde van het jaar worden de eigenschappen van bewerkingen onder de loep genomen.
Leerlingen onderzoeken eerst wat er precies kan en mag met concrete getallen en voeren daarna een veralgemening door.
Ze formuleren de eigenschappen met letters, waarbij deze letters staan voor eender welk getal uit een afgebakende verzameling.
Deze eigenschappen worden dan toegepast enerzijds bij handige rekenmethodes (net zoals in de lagere school) en anderzijds om te letterrekenen.
Dit laatste is helemaal nieuw.

\subsection{Mind the gap}

Op het eerste gezicht lijkt het alsof twaalfjarigen helemaal klaar zijn voor de stap naar het secundair onderwijs.
Ze leerden alle bewerkingen met natuurlijke getallen uitvoeren en ze leerden ook de belangrijkste eigenschappen van deze bewerkingen kennen en gebruiken om handig te rekenen.
Daarnaast leerden ze ook nog (een aantal) bewerkingen met breuken uit te voeren.

Toch zullen leraren uit de eerste graad van het secundair onderwijs beamen dat er wel een `gap' te bespeuren valt.
Het grootste probleem is wellicht dat de leerlingen al die verschillende regeltjes (methodes en eigenschappen) uit de lagere school ergens een plek in hun hoofd hebben gegeven.
Ze hebben echter te weinig zicht op de onderliggende structuren (welk statuut heeft een eigenschap en binnen welke verzameling mag ze worden toegepast?) en de onderlinge samenhang (bv. splitsen en verdelen en factoren samennemen zijn allebei distributiviteit, maar dan omgekeerd).
Ze hebben niet per se een goed overzicht van wat er allemaal bestaat en we kunnen er dus niet van uitgaan dat leerlingen in het eerste middelbaar vlot en flexibel alle aangeleerde regeltjes kunnen inzetten wanneer ze een gevarieerde set aan oefeningen voor zich krijgen.
Daarnaast of hiermee samenhangend zijn ook een meer formeel wiskundig taalgebruik en correcte wiskundige notaties voor veel leerlingen een plotse verandering, wat als een, vaak onderschatte, drempel wordt ervaren.

\subsection{Hoe bouw je de brug?}
Om een vlotte overgang naar de volgende stap in abstrahering naar o.a. letterrekenen te kunnen maken, is het goed om op deze moeilijkheden te focussen.

\subsubsection*{Het bos en de bomen}
Het is verleidelijk om een oplossingsmethode al snel te veralgemenen door deze met behulp van letters te noteren.
Wanneer een wiskundeleraar het principe van commutativiteit bespreekt, zal hij al snel \(a + b = b + a\) op het bord schrijven.
Ook in de leerwerkboeken vinden we deze notatie al in de eerste hoofdstukken van het eerste jaar secundair onderwijs terug.
Als je bedenkt dat meerdere leerlingen in je klas nog aan het worstelen zijn met het herkennen van opgaven waarin ze commutativiteit, of `van plaats wisselen', mogen toepassen, dan is het wellicht zinvoller om hen eerst een overzicht te geven met behulp van concrete getallen.
In \autoref{fig:voorbeeldeig} vind je een voorbeeld van zo'n overzicht terug.
Het is beperkt tot de vermenigvuldiging, bij wijze van voorbeeld.
Een uitgebreider overzicht met alle bewerkingen kan online bij deze loep op \url{uitwiskeling.be} geraadpleegd worden.

\end{multicols}
\afbeelding[14cm]{img/loepMichele/figuurm12.png}{Overzicht met voorbeelden van de regels voor producten}{fig:voorbeeldeig}
\begin{multicols}{2}

Via zo'n overzicht krijgen leerlingen ook meer grip op de onderliggende structuren en de onderlinge samenhang van de bewerkingen en hun eigenschappen.
Er is in de eerste plaats voor elke bewerking een standaardmethode.
Daarnaast bestudeer je de verschillende eigenschappen, zijnde `van plaats wisselen', `schakelen' en `splitsen en verdelen'.
Vooral voor de optelling en de vermenigvuldiging leveren die al heel wat rekenvoordelen op.
Bij de aftrekking en de deling is het nodig om alert te zijn, aangezien de eigenschappen slechts gedeeltelijk toepasbaar zijn. Denk bijvoorbeeld aan de rechts-distributiviteit van de deling.
Tenslotte zijn er ook nog andere rekenregels die rekenvoordelen opleveren.
Die hebben meestal met compenseren te maken, steunend op de omgekeerde bewerking.
Wanneer alle bewerkingen samenkomen, is het nodig om je aan de voorrangsregels te houden, en de opgaven op te lossen rekening houdend met de volgorde van bewerkingen.

\subsubsection*{Correct wiskundig taalgebruik en wiskundige notaties}
Hoewel leraren in de basisschool zeker begrippen zoals termen en factoren gebruiken, zijn veel leerlingen het niet gewend om deze taal zelf te hanteren en ook niet om ze zonder concreet voorbeeld te interpreteren.
Om de overgang vlot te maken, kun je als leraar in het eerste jaar van het secundair onderwijs best een aantal voorbeelden, zoals die uit het overzicht hierboven, paraat hebben.

Ook de formele benamingen van de wiskundige eigenschappen commutativiteit, associativiteit en distributiviteit, gebruiken we best nog een voldoende lange tijd samen met de eerder gekende informele omschrijvingen `van plaats wisselen', `schakelen' en `splitsen en verdelen'.

Tenslotte is het ook goed om te weten dat wiskundemethodes voor het basisonderwijs de afgelopen jaren op zoek zijn gegaan naar een soort van algoritmes of notaties waarmee ze pogen de gewenste denkwijze te visualiseren.
In het derde leerjaar wordt in de eerste les over distributiviteit (splitsen en verdelen als standaardmethode voor het vermenigvuldigen buiten de tafels)  bijvoorbeeld de notatie uit \autoref{fig:splitsbeentjes} met splitsbeentjes gebruikt.

\end{multicols}
\afbeelding[10cm]{img/loepMichele/figuurm13.png}{Visualisatie van distributiviteit in het basisonderwijs (Uit: \emph{Wiskanjers 3}, blok 3, les 19)}{fig:splitsbeentjes}
\begin{multicols}{2}
Hoewel deze manier van noteren uiteraard met de beste bedoelingen door deze methodes en leraren wordt gebruikt, leren leerlingen niet wat de distributiviteit eigenlijk doet en leren ze die ook niet wiskundig noteren.
Dat geeft problemen wanneer je later diezelfde distributiviteit bijvoorbeeld wilt gebruiken bij handig rekenen.
In hetzelfde werkboek wordt in het vijfde leerjaar compenseren als handige methode aangereikt (zie \autoref{fig:handycommu}). Slechts zelden leggen leerlingen de link tussen beide.

\end{multicols}
\afbeelding[14cm]{img/loepMichele/figuurm14.png}{Visualisatie van compenseren in het basisonderwijs (Uit: \emph{Wiskanjers 5}, blok 3, les 13)}{fig:handycommu}
\begin{multicols}{2}



Leraren in de eerste graad van het secundair onderwijs kunnen de overgang vergemakkelijken door net deze verbanden te expliciteren.
Wanneer je vertrekt van de opgave \(13 \times 6\), kan je de notatie met de splitsbeentjes geleidelijk aan vervangen door de correcte wiskundige notatie.
De eerste tussenstap is hierbij de meest cruciale, aangezien deze duidelijk maakt wat de distributiviteit precies inhoudt.
\[
13 \times 6 = (10+3)\times 6 = (10 \times 6) + (3 \times 6) = 60 + 18 = 78
\]

\section{Getallenas en assenstelsel}\label{getallen}
\subsection{De getallenas in het basisonderwijs}
We staan eerst even stil bij de benaming.
Het handboek dat we hieronder als illustratie gebruiken, spreekt over de `getallenas'.
In de (huidige) eindtermen voor de basisschool (die in een proces van herziening zitten) wordt de naam `getallenlijn' gebruikt.
Soms horen we ook dat beide namen gebruikt worden, met een verschil in betekenis.
Dan is de getallenlijn een rij getallen zonder ijk.
Enkel de ordening van de getallen is dan van belang en niet de grootte.
In deze loep gebruiken we de naam \emph{getallenas}.

In de basisschool zien we de getallenas geregeld terugkomen.
Leerlingen moeten getallen op een getallenas plaatsen of, omgekeerd, aangeven met welk getal een punt op een getallenas overeenkomt.
Leerlingen moeten dit kunnen met natuurlijke getallen, met positieve breuken en decimale getallen, en er wordt ook een aanzet gegeven om eenvoudige negatieve getallen op de getallenas te situeren in concrete contexten (bijvoorbeeld: temperatuur, lift).
In de meeste gevallen krijgen de leerlingen een getallenas waar de \(0\) en \(1\) al vooraf op aangeduid zijn.
Als er breuken voorgesteld moeten worden, is er vaak ook al een passende onderverdeling van de eenheid gemaakt, zoals in \autoref{fig:image1}.
\end{multicols}
\afbeelding[14cm]{img/loepJohan/image1.png}{Een reeds geijkte getallenas (Uit: \emph{Wiskanjers 6}, blok 1, les 2)}{fig:image1}
\begin{multicols}{2}
De vooraf aangeduide getallen zijn niet altijd \(0\) en \(1\), waardoor leerlingen dus soms een minder klassiek deel van de getallenas te zien krijgen.
In \autoref{fig:image2} gaat het over een fijnere onderverdeling van de getallenas, maar het kan ook gaan over het voorstellen van heel grote getallen.
\end{multicols}
\afbeelding[14cm]{img/loepJohan/image2.png}{Een fijner geijkte getallenas (Uit: \emph{Wiskanjers 6}, blok 1, S3)}{fig:image2}
\afbeelding[14cm]{img/loepJohan/image3.png}{Een ongeijkte getallenas (Uit: \emph{Wiskanjers 6}, blok 1, les 2)}{fig:image3}

\begin{multicols}{2}


\end{multicols}
\afbeelding[14cm]{img/loepJohan/image4.png}{Sprongen op een getallenas (Uit: \emph{Wiskanjers 6}, blok 1, les 5)}{fig:image4}
\begin{multicols}{2}


\end{multicols}
\afbeelding[14cm]{img/loepJohan/image5.png}{Een getallenas in context (Uit: \emph{Wiskanjers 6}, oplossingen, blok 5, les 3)}{fig:image5}

\newpage

\begin{multicols}{2}

Oefeningen waarbij leerlingen zelf een gepaste ijk moeten kiezen, zoals in \autoref{fig:image3}, zijn een uitzondering. 

De getallenas wordt in verband gebracht met de ordening van getallen.
Ze wordt ook gebruikt bij bepaalde strategieën voor het uitvoeren van bewerkingen.
In \autoref{fig:image4} zie je een opgave (met oplossing) waarbij de vermenigvuldiging van een breuk met een natuurlijk getal weergegeven wordt met sprongen.
Voor delen kun je ze gebruiken voor de verhoudingsdeling, bijvoorbeeld om te bepalen hoeveel groepjes van \(3\) er gaan in \(15\).
Een (lege) getallenas (waarop je de groottes niet echt afmeet maar schat) leent zich goed om optellingen op te lossen met de rijgmethode, waarbij we enkel de tweede term splitsen, zoals bijvoorbeeld in \(23 + 48 = 23 + 40 + 8\).
Dit is een betere werkwijze (minder geheugenbelasting, minder kans op fouten) dan de splitsprocedure: \(23 + 48 = (20 + 40) + (3 + 8)\).
De getallenas is echter niet het enige hulpmiddel: voor bepaalde andere strategieën is het bijvoorbeeld nuttiger om materiaal in te zetten dat de \(10\)-structuur van ons talstelsel ondersteunt.

We hebben in het handboek dat we geconsulteerd hebben geen contextopgaven gevonden waarin leerlingen gebruik moesten maken van een getallenas.

Tot slot, de modeloplossing uit \autoref{fig:image5} is opmerkelijk als we ze bekijken met het begrip `getallenas' in het achterhoofd.
De werkschets is eigenlijk niets anders dan een verticaal geplaatste getallenas, dit keer wel in een context, waarbij leerlingen zelf de ijk moeten kiezen en waarbij een visuele voorstelling van de bewerkingen de leerlingen moet helpen om het vraagstuk op te lossen.


\subsection{De getallenas in het secundair onderwijs}\label{de-getallenas-in-het-secundair-onderwijs}

In de eindtermen wordt de getallenas genoemd bij het ordenen van getallen.
Het wekt dan ook weinig verwondering dat we bij het ordenen van getallen de getallenas aantreffen in de handboeken.
Negatieve getallen komen in de eerste graad van het secundair onderwijs nu wel systematisch aan bod en dat zorgt natuurlijk ook voor een aanpassing van de getallenas, die nu in beide richtingen oneindig doorloopt.
Vaak wordt het verband gelegd met de schaal op een thermometer en met een tijdslijn uit de geschiedenis (zie \autoref{fig:image6}). 
\end{multicols}
\afbeelding[14cm]{img/loepJohan/image6.png}{Thermometer als getallenas (Uit: \emph{Delta Nova 1A}, p. 221)}{fig:image6}
\begin{multicols}{2}
Begrippen als tegengestelde en absolute waarde, die hun oorsprong vinden bij het werken met negatieve getallen, worden geduid op de getallenas, zoals te zien is in \autoref{fig:image7}.
Verder lijkt het allemaal nogal op wat leerlingen reeds in de basisschool geleerd hebben: bewerkingen worden weinig of niet geïnterpreteerd met behulp van de getallenlijn, de leerlingen moeten zelden of nooit zelf een ijk kiezen en oefeningen waarin een getallenlijn in een context voorkomt zijn zeldzaam of zelfs helemaal afwezig. % todo: deze paragraaf niet splitsen bij voorkeur van Johan, wel indien beter met figuren
\end{multicols}

\afbeelding[14cm]{img/loepJohan/image7.png}{Tegengestelde en absolute waarde op een getallenas (Uit: \emph{Van basis tot limiet 1}, leerboek, p. 59)}{fig:image7}
\begin{multicols}{2}

\subsection{Coördinaten en assenstelsels in het secundair onderwijs}\label{couxf6rdinaten-en-assenstelsels-in-het-secundair-onderwijs}

De getallenas is een manier om getallen meetkundig voor te stellen in één dimensie.
Je kunt het ook omgekeerd bekijken: door van een rechte een getallenas te maken, kun je de positie van een punt beschrijven met een getal.
Dit komt overeen met wat we doen met coördinaten ten opzichte van een assenstelsel: we beschrijven de positie van een punt in het vlak door middel van twee coördinaatgetallen, waarbij we gebruik maken van een assenstelsel dat in feite een combinatie is van twee getallenassen.
In de eerste graad van het secundair onderwijs wordt hiervoor de aanzet gegeven.

\end{multicols}
\afbeelding[14cm]{img/loepJohan/image8.png}{Uitbreiding naar negatieve coördinaten (Uit: \emph{Delta Nova 1A}, p. 224)}{fig:image8}
\begin{multicols}{2}

In beide geconsulteerde handboeken zagen we dat coördinaten gerelateerd worden aan vormen van plaatsbepaling die leerlingen al kennen.
Het kan daarbij gaan over kennis uit het dagelijkse leven (bijvoorbeeld de velden op een schaakbord of bij het spel zeeslag) of om plaatsbepaling die leerlingen in het basisonderwijs al ontmoet hebben, maar dan buiten het vak rekenen/wiskunde.
Denk bij dit laatste bijvoorbeeld aan het gebruik van een letter en een getal om een vak op een stadsplan aan te duiden.

Soms wordt aanvankelijk enkel met gehele of met positieve coördinaatgetallen gewerkt en volgt de uitbreiding tot negatieve (zie \autoref{fig:image8}) of rationale coördinaatgetallen later.
We zien in \autoref{fig:image8} ook dat aan het assenstelsel soms een rooster gekoppeld wordt, waardoor punten met gehele coördinaatgetallen duidelijk zichtbaar worden.
In \autoref{fig:image8} zien we dat de coördinaatgetallen betekenis krijgen via stappen naar rechts/links en naar boven/onder.
Bij de oefeningen zien we dan oefeningen waarbij leerlingen vertrekken vanuit een gegeven punt en de coördinaten moeten bepalen van een ander punt dat je vindt door horizontaal of verticaal een aantal stappen te zetten.
In andere gevallen worden de coördinaatgetallen in de eerste plaats geïnterpreteerd via projectie van een punt op beide assen in plaats van (of in combinatie met) de interpretatie met de stapjes.

\subsection{Mind the gap en maak de brug}

\subsubsection*{Assenstelsels en coördinaten}

Bij assenstelsels en coördinaten vinden we het goed dat de handboeken aansluiten bij wat leerlingen op het vlak van plaatsbepaling in het basisonderwijs al hebben leren kennen.
Toch zien we hier een hobbel(tje) in de leerlijn.
Bij het lezen van een stadsplan hebben leerlingen geleerd om een vak op de kaart aan te duiden met een letter en een getal, zoals in \autoref{fig:image9}.
Als iemand \(4B\) gebruikt om een vak aan te duiden i.p.v. \(B4\) is er geen probleem. Bij coördinaten in de wiskunde gebruiken we twee getallen en dan is de volgorde juist wel van belang.

\afbeelding{img/loepJohan/image9.png}{Coördinaten op een stadsplan (Uit: \emph{Van basis tot limiet 1}, leerboek, p. 61)}{fig:image9}

Er is nog een tweede verschil, en dat is fundamenteler.
Op het stadsplan gebruiken we de `coördinaten' om een heel vak aan te geven (met daarin oneindig veel punten), terwijl de coördinaten in de wiskunde één punt bepalen, zoals bijvoorbeeld de roosterpunten \(A\), \(B\), \(C\) en \(D\) in \autoref{fig:image8}.
Voor een goed begrip is het belangrijk om leerlingen te wijzen op dit verschil.
Het koppelen van een rooster aan een assenstelsel is op zich een goed idee, maar als je niet oplet kan er hier verwarring optreden.
Voor de leraar: we maken hier de overgang van discreet (stadsplan) naar continu (meetkunde).

\subsubsection*{Getallenas}

Wat ons opvalt omtrent de getallenas, is dat er in feite weinig nieuws aan bod komt in de eerste graad van het secundair onderwijs (met uitzondering van het voorstellen van negatieve getallen en meer gevorderd werken met breuken).
Net als in het basisonderwijs zien we bijvoorbeeld dat leerlingen in de eerste graad van het secundair onderwijs zelden of nooit zelf een gepaste ijk moeten kiezen.
Nochtans zullen ze in de hogere jaren van het secundair onderwijs geregeld zelf een geschikte keuze moeten maken voor de eenheden op de assen van een assenstelsel: in functie van de context, om relevante aspecten van het verloop van een functie te kunnen zien, om het tekenvenster te kunnen instellen, \ldots{}
Denkend aan het verband dat we gelegd hebben tussen de getallenas en assenstelsels, zien we het daarom als een zinvolle stap in de leerlijn als leerlingen in de eerste graad ook al eens een gepaste ijk op een getallenlijn moeten kiezen.

Een voorbeeld van een contextloze opgave waarin dit gevraagd wordt, is de volgende: plaats \(\frac{3}{4}\), \(\frac{2}{3}\) en \(\frac{7}{4}\) op de getallenas.
Leerlingen moeten er dan enerzijds rekening mee houden dat de ijk `deelbaar' moet zijn door \(4\) en door \(3\) (dus, bijvoorbeeld, \(12\) ruitjes), en anderzijds dat ook \(2\) (of iets minder) op de getallenas moet passen.


Voor een (moeilijke!) contextopgave spelen we leentjebuur bij een opgave die we in een handboek van het basisonderwijs vonden (zie \autoref{fig:image10}).

\end{multicols}
\afbeelding{img/loepJohan/image10.png}{Contextopgave voor zelf ijken van getallenas (Uit: \emph{Wiskanjers 6}, blok 2, les 6)}{fig:image10}
\begin{multicols}{2}
Ze werd daar gebruikt om leerlingen uit het basisonderwijs (benaderend) te laten rekenen met grote getallen. Aan leerlingen uit de eerste graad van het secundair onderwijs geven we een andere opdracht: stel alle afstanden voor op een getallenas.
Je moet je dan bijvoorbeeld afvragen welke schaal je zult gebruiken (bij een A4-blad in `landscape' kun je \(\SI{1}{\centi\metre}\) laten overeenkomen met \(200\) miljoen kilometer zodat Neptunus nog op je blad past) en je kunt je afvragen of de \(0\) wel zichtbaar moet zijn (niet per se, maar best wel, denken we: daar situeert de zon zich. Het vergt ook nauwelijks bijkomende plaats omdat de zon en Mercurius zo dicht bij elkaar liggen).
Als je de oefening gemaakt hebt, zul je vaststellen dat de illustratie in de opgave niet bepaald op schaal is\ldots{}
Deze opgave is uitdagend omwille van de grote getallen, maar de overwegingen die we gemaakt hebben, zijn typisch voor wat aan bod komt bij het kiezen van geschikte eenheden op de assen wanneer we in een context een functie moeten voorstellen.
Daar komt het natuurlijk op aan en je kunt de leerlingen dit ook laten inoefenen in opgaven met kleinere getallen.

\subsubsection*{Nogmaals assenstelsels en coördinaten}
We hebben assenstelsels en coördinaten hierboven verantwoord vanuit plaatsbepaling: we leggen de positie van een punt vast door zijn coördinaatgetallen te geven ten opzichte van een assenstelsel.
Dit is het startpunt voor het gebruik van assenstelsels in de analytische meetkunde.
Assenstelsels worden echter ook in andere deeldomeinen van de wiskunde gebruikt, namelijk in de analyse (grafisch voorstellen van functies) en in de statistiek (lijndiagram, spreidingsdiagram, \ldots).

In de (elementaire) analytische meetkunde staan de assen typisch loodrecht op elkaar en zijn de eenheden even groot.
Bij het grafisch voorstellen van functies staan de assen wel nog loodrecht op elkaar, maar is het niet altijd mogelijk en is het ook niet altijd zinvol om de eenheden even groot te nemen.
Neem bijvoorbeeld de functie \(f\) die de prijs \(p\) van drukwerk (in \(\euro\)) geeft in termen van het aantal gewenste afdrukken \(a\), met als vergelijking \(p = f(a) = 5 + \num{0.03}a\).
Omdat de uitgezette getallen een verschillende grootteorde hebben (bijvoorbeeld: \(1000\) afdrukken kosten \(\EUR{35}\)), zijn gelijke eenheden niet handig.
Maar er is ook een diepere reden om hier niet voor gelijke eenheden te kiezen: op de verticale as zetten we bedragen uit en op de horizontale as komen aantallen te staan.
Met andere woorden: bij het grafisch voorstellen van functies in een context zijn de veranderlijken op beide assen (doorgaans) niet meer `gelijksoortig'.
Dat is ook zo bij de lijndiagrammen uit de statistiek (en bij spreidingsdiagrammen, die in de tweede graad op het programma staan).
In de analytische meetkunde ligt dat anders: daar zetten we op beide assen afstanden uit.

\subsubsection*{Een uitweiding naar de tweede en de derde graad}

In de eerste graad van het secundair onderwijs blijven de gevolgen van al dan niet gelijksoortig zijn van de twee veranderlijken beperkt tot het al dan niet kiezen van gelijke eenheden op de assen.
Als je verder gaat in de wiskunde, dan zijn de gevolgen groter.
Denk bijvoorbeeld aan de afstandsformule:
\[
d(a,\ b) = \sqrt{\left( x_{b} - x_{a} \right)^{2} + \left( y_{b} - y_{a} \right)^{2}}
\]
Als \(x\) en \(y\) niet gelijksoortig zijn, dan heeft de som van de twee termen onder het wortelteken geen betekenis.
Bij de assenstelsels die we in de analyse en de statistiek gebruiken, kun je wel nog horizontale en verticale `afstanden' bepalen, maar `schuine afstanden' zijn niet zinvol.
Een concreet gevolg hiervan zie je bijvoorbeeld bij de kleinstekwadratenrechte (trendlijn) in de tweede graad: daar wordt gewerkt met verticale afstanden tussen punten en de rechte.
Wat we hier vaststellen voor de afstandsformule geldt voor heel veel van wat we kennen uit de analytische meetkunde: de vergelijking van een cirkel, de betrekking tussen richtingscoëfficiënten van loodrecht op elkaar staande rechten, de interpretatie van de richtingscoëfficiënt van een rechte als tangens van de hoek met de horizontale as, \ldots{}

Hier past tot slot nog een nuancering: het verschil dat we hierboven besproken hebben, is niet altijd even helder.
Nu en dan zijn beide veranderlijken in een functie of in een spreidingsdiagram wél gelijksoortig.
Denk bijvoorbeeld aan de functie die het verband tussen lengte en breedte beschrijft van een rechthoek met een vaste oppervlakte.
Een ander voorbeeld is het gebruik van functies in de analytische meetkunde om krommen te beschrijven.
En, tot slot: wanneer je functies los van enige context bestudeert, dan hebben de veranderlijken \(x\) en \(y\) sowieso geen concrete betekenis.

\section{Meetkunde}\label{meetkunde}
Doorheen de voorbije curriculumhervormingen heeft meetkunde sterk aan ruimte ingeboet en is op een of andere manier in het verdomhoekje beland.
Verschillende meetkundige onderwerpen hebben plaats moeten ruimen ten voordele van andere wiskundeonderdelen zoals dataverwerking, statistiek en computationeel denken.
Daardoor is een goed opgebouwde leerlijn in de leerplannen en handboeken nog moeilijk te vinden.
Dat is een gemiste kans want meetkunde is een rijk onderwerp waarmee leerlingen kunnen oefenen in het opbouwen van correcte redeneringen.
En net oefenen met deze manier van denken zou het wiskundig redeneren, waaronder inzichten voor het computationeel denken, zeker ten goede komen.

\subsection{Leerstof in het basisonderwijs}
In het basisonderwijs verkennen en onderzoeken de leerlingen de meest courante meetkundige objecten.
Leerlingen lossen oefeningen op door iets te tekenen, aan te kruisen, te meten en in te vullen.
Een typische opgave vind je in \autoref{fig:diagonalen}.

Op basis van dit soort opgaven doen leerlingen vaststellingen i.v.m. meetkundige figuren en hun eigenschappen.
Een aantal onderliggende meetkundige inzichten -- die zelden worden geëxpliciteerd -- spelen in deze opgaven een rol:

\begin{itemize}
\item Inzicht in het inclusieprincipe of het concept deelverzameling.
Voor meetkundige objecten betekent dit dat je begrijpt en beseft dat een vierkant ook een rechthoek, een ruit \ldots{} is.
\item Op basis van één voorbeeldfiguur kun je niet zomaar een algemene eigenschap afleiden, zoals `de diagonalen van een rechthoek snijden elkaar middendoor'.
\item De wiskundige verwoording volgt een aantal regels, die in het dagelijkse leven anders zijn.
Bv. `de vierhoeken met één paar evenwijdige zijden omvatten de vierhoeken met twee paar evenwijdige zijden'.
\end{itemize}
\end{multicols}
\afbeelding[14cm]{img/loepEls/image1.png}{Diagonalen (Uit: \emph{Wiskidz 6}, blok 2, les 5)}{fig:diagonalen} 
\begin{multicols}{2}


Soms zijn er in de opgave elementen die hier iets zichtbaar van maken: `Noteer de meest \emph{passende} naam'.
Al is het niet waarschijnlijk dat leerlingen de betekenis van het woordje `passende' zelf doorzien.
Verschillende opgaven bevatten echter elementen die dit soort inzichten in de weg staan.
Zo wekt de zin `een vierhoek waarbij de diagonalen loodrecht op elkaar staan, is een \emph{ruit of een vierkant}' (\autoref{fig:diagonalen}) de indruk dat een vierkant geen ruit zou zijn. Bovendien zijn er nog  andere vierhoeken met loodrechte diagonalen die geen ruiten zijn. Een correcte formulering zou zijn: `Een parallellogram waarbij de diagonalen loodrecht op elkaar staan, is een ruit.'

Een alerte leraar vangt deze lacunes op door goede denkvragen te stellen en aan te dringen op correcte verwoordingen: `Is dat voor alle rechthoeken zo of alleen voor dit voorbeeld?', `Kun je een trapezium tekenen waarvan de diagonalen wél gelijk zijn?' en `Is het dan niet voldoende dat we zeggen dat een parallellogram waarbij de diagonalen loodrecht op elkaar staan een ruit is?'\ldots

\subsection{Leerstof in het secundair onderwijs}
De meetkundige objecten die in het secundair onderwijs aan bod komen, zijn in grote mate dezelfde.
Hooguit een aantal specifiekere objecten zoals een prisma, een halfrechte, een bissectrice\ldots{} zijn nieuw voor de leerlingen.
In handboeken zien we dat het formele aspect meer aandacht krijgt: het formuleren van definities, het correct benoemen van hoekpunten met letters\ldots{}
Ook wordt er meer aandacht aan het inclusieprincipe besteed.
In de bovenstaande opgave \autoref{fig:vennmeetkunde} wordt met behulp van venndiagrammen duidelijk gemaakt dat een vierkant ook een ruit of rechthoek is.
In \autoref{fig:uitsprakenmeet} zie je dat eigenschappen formeler met implicaties en kwantoren worden geformuleerd.

De theorie van Van Hiele is een hulpmiddel dat aan het meetkundeonderwijs richting geeft.
Deze theorie gaat er vanuit dat meetkunde zich op verschillende denkniveaus kan afspelen.
Een overzicht vind je terug in de onderstaande tabel \autoref{tab:vanhiele} (Thaels,
Eggermont, \& Janssens, 2001).

\end{multicols}
\afbeelding[14cm]{img/loepEls/image2.png}{Meetkundig venndiagram (Uit: \emph{Delta Nova 1}, H7, p. 218)}{fig:vennmeetkunde}

\afbeelding[14cm]{img/loepEls/image3.png}{Meetkundige uitspraken (Uit: \emph{Delta Nova 1}, H7, p. 218)}{fig:uitsprakenmeet}
\begin{multicols}{2}

\end{multicols}

\begin{tabel}
	\begin{tabular}{|m{2.5cm}|m{9cm}|m{3cm}|}
		\hline
		\tableheader{Niveau} & \tableheader{Omschrijving} & \tableheader{Uitspraken van leerlingen}\\ 
		\hline\rowcolor{white}
		1 Visualisatie  & Leerlingen herkennen meetkundige figuren `op het zicht' door ze te vergelijken met een prototype.
        Eigenschappen van figuren worden op dit niveau niet onderkend.
        Een vierkant dat op zijn punt gezet is, noemen ze vaak hardnekkig een ruit, ook al draai je voor hun ogen een vierkant op een punt. & Deze figuur is een rechthoek omdat ze lijkt op een deur. \\ 
		\hline\rowcolor{white}
		2 Analyse & Leerlingen zien figuren als een verzameling van eigenschappen.
        Ze kunnen eigenschappen herkennen en benoemen maar zien er geen verband tussen.
        Als een leerling een figuur beschrijft, zal hij alle eigenschappen benoemen zonder onderscheid tussen noodzakelijke en niet noodzakelijke eigenschappen.
        De leerling heeft geen besef van wat `voldoende' eigenschappen zijn om een figuur te beschrijven.
        Ze zien ook geen verband tussen verschillende figuren. & Een vierkant is geen rechthoek.\\ 
		\hline\rowcolor{\rubriekkleur!30}
		3 Abstrahering  & Leerlingen zien verbanden tussen eigenschappen en tussen figuren onderling.
        Ze kunnen \textbf{betekenisvolle definities} geven en \emph{informele argumenten} formuleren om hun redenering te staven.
        Ze zien een verband tussen de verschillende \textbf{klassen van figuren}.
        Ze kunnen figuren hiërarchisch ordenen.
        Ze zijn in staat eenvoudige \textbf{`als.. dan\ldots'- redeneringen} op te bouwen.
        De rol en betekenis van een formele deductie begrijpen ze echter nog niet. & Een vierkant heeft alle eigenschappen van een rechthoek en is dus een rechthoek.\\ 
		\hline\rowcolor{white}
		4 Ordening  & Leerlingen kunnen bewijzen opstellen.
        Ze begrijpen de rol van ongedefinieerde begrippen, definities, axioma's en stellingen.
        Ze zien in dat om de juistheid van een stelling na te gaan, het niet volstaat om zeer veel voorbeelden te controleren, maar dat je door redenering een `bewijs' moet geven.
        Ze vatten de bedoeling van een strenger en logischer systeem waarin de eigenschappen van de figuren passen.
        Ze kunnen werken met abstracte uitspraken en trekken besluiten op basis van logische redeneringen en niet op basis van intuïtie. & Ik kan bewijzen dat de diagonalen van een vierhoek elkaar in het midden snijden, als de vierhoek een parallellogram is.\\ 
		\hline\rowcolor{white}
		5 Wiskundig systeem  & Leerlingen begrijpen de formele aspecten van deductie zoals het opstellen en vergelijken van wiskundige systemen.
        Ze kunnen indirecte bewijzen en bewijzen door contrapositie hanteren en ze zijn klaar voor niet-euclidische meetkunde. & `Een vierkant is een rechthoek' is hetzelfde als `een vierhoek die geen rechthoek is, is ook geen vierkant'.\\
		\hline
	\end{tabular}
	\caption{Van Hiele theorie}\label{tab:vanhiele}
\end{tabel}

\begin{multicols}{2}
In de derde graad van het basisonderwijs en de eerste graad van het secundair onderwijs zijn leerlingen bezig met het verwerven van niveau 3.
Vaak wordt er gedacht dat leerlingen voor een bepaald niveau `rijp' moeten zijn.
Van Hiele kent echter een belangrijke rol toe aan de leraar.
De overgang van het ene niveau naar een volgend hangt meer van goed onderwijs af dan van de leeftijd of maturiteit van de leerling.
In de volgende paragrafen zoomen we verder in op `betekenisvolle definities', redenering met `informele argumenten', `classificatie van meetkundige objecten' en `als \ldots{} dan \ldots{} redeneringen' uit niveau 3.

\subsection{Betekenisvolle definities}\label{betekenisvolle-definities}

\subsubsection*{Mind the gap}
De overgang van niveau 2 naar niveau 3 gebeurt in de meeste gevallen eerder impliciet.
Leerlingen en leraren zijn zich er niet altijd van bewust dat op een nieuw niveau wordt geredeneerd.
Ook handboeken besteden aan deze overgang weinig aandacht.
De formulering van de opgaven werkt deze abstrahering niet altijd in de hand.
De volgende twee voorbeelden illustreren dit.
\end{multicols}
\afbeelding[14cm]{img/loepEls/image4.jpeg}{Meetkundefiguren herkennen (Uit: \emph{Delta Nova 1}, H2, p. 113)}{fig:herkennen}
\begin{multicols}{2}

De opgave uit \autoref{fig:herkennen} past beter bij niveau 1 waar leerlingen figuren `herkennen' dan bij het niveau `abstrahering'.
Wanneer er veel van dit soort opgaven aanwezig zijn, leidt dit de leerlingen weg van het volgende niveau.
De opgave in \autoref{fig:verschiltrap} toont hoe vraagstelling de groei naar een volgend niveau in de weg kan staan.

\end{multicols}
\afbeelding[14cm]{img/loepEls/image5.png}{Verschil trapezium en parallellogram (Uit: \emph{Wiskidz 6}, blok 2, les 5)}{fig:verschiltrap}
\begin{multicols}{2}

Door te vragen naar het \textbf{verschil} tussen een trapezium en een parallellogram, wek je impliciet de indruk dat dit twee onderscheiden begrippen zijn die naast elkaar bestaan.
Terwijl het parallellogram een speciaal trapezium is waardoor het net alle eigenschappen van het trapezium automatisch bezit.

\subsubsection*{Hoe bouw je de brug?}
Als leraar ondersteun je leerlingen door hun ervaringen te structuren zodat ze gemakkelijker de stap naar het volgend niveau kunnen zetten.
Daarbij is het belangrijk te beseffen dat elk niveau zijn eigen taal heeft.
Typisch voor niveau 3 is het gebruik van definities om een begrip eenduidig vast te leggen.
Om te vermijden dat leerlingen `definities zomaar van buiten blokken en aframmelen' is het zinvol om het \textbf{concept} definitie te verduidelijken.
In de volgende lesactiviteit geven we een voorbeeld hoe je daaraan kunt werken.
\end{multicols}

\begin{lesactiviteit}{Definities}
Meetkundige objecten benoemen we met een (verzamel)naam zoals vierkant, kubus, trapezium\ldots{}
We leggen deze begrippen vast met een omschrijving zodat je ze niet met andere begrippen kunt verwarren.
We noemen dit een definitie.
Bv.:
\begin{citaat}
    Een parallellogram is een \underline{vierhoek} met twee paar evenwijdige zijden.
\end{citaat}

Een definitie is sober en bevat het minimum aan eigenschappen of kenmerken om het begrip vast te leggen.
De onderstaande omschrijving voor een parallellogram gebruiken we niet als definitie.
Ze bevat meer informatie dan nodig om het begrip van andere begrippen te
onderscheiden.
\begin{citaat}
    Een parallellogram is een \underline{vierhoek} met twee paar evenwijdige zijden die even lang zijn, waarvan de overstaande hoeken gelijk zijn, met diagonalen die elkaar middendoor snijden.
\end{citaat}

Soms zijn er meerdere varianten van een definitie mogelijk.
Bijvoorbeeld:
\begin{itemize}
\item Een vierkant is een \underline{vierhoek} met vier gelijke zijden en vier gelijke hoeken.
\item Een vierkant is een \underline{rechthoek} met vier gelijke zijden.
\item Een vierkant is een \underline{ruit} met vier gelijke hoeken.
\end{itemize}

Sommige begrippen, zoals een punt, een rechte, een vlakke figuur\ldots{} kunnen we (nog) niet definiëren, maar gebruiken we vanuit intuïtief inzicht.

Om definities te maken is het handig dat je een algemener begrip gebruikt.
In de bovenstaande definities onderstreepten we dat begrip: een vierhoek, rechthoek, ruit.
Het volgende voorbeeld maakt duidelijk waarom een algemener begrip nodig is:
\begin{itemize}
\item `Een diagonaal verbindt twee niet opeenvolgende hoekpunten.'\\
  Deze omschrijving bevat geen bovenliggend begrip.
  Daarom sluit ze de onderstaande kromme niet uit als diagonaal.
  \afbeelding[6cm]{img/loepEls/vreemdediagonaal.png}{}{}

\item De omschrijving met het bovenliggende begrip doet dat wel: `Een diagonaal is een \underline{lijnstuk} dat twee niet opeenvolgende hoekpunten verbindt.
\end{itemize}

{\color{\rubriekkleur}\section*{Opdrachten}}
    
\begin{vraagenantwoord}
\vraag{Waarom zijn volgende zinnen geen goede definities?}
\begin{vraagenantwoord}
\vraag{In een gelijkzijdige driehoek zijn alle zijden gelijk.}
    \antwoord{Het bovenliggende begrip ontbreekt.}
    \vraag{Een vierkant is een vierhoek met vier gelijke zijden.}
    \antwoord{Deze uitspraak is onvolledig.}
    \vraag {Een ruit is een vierhoek met vier gelijke zijden en vier gelijke hoeken waarvan de overstaande hoeken gelijk zijn.}
    \antwoord{Deze uitspraak bevat foute informatie; een ruit hoeft geen vier gelijke hoeken te hebben.}
   \vraag{Een symmetrieas spiegelt de figuur op zichzelf.}
    \antwoord{Het bovenliggende begrip ontbreekt.}
    \vraag{Een regelmatige vijfhoek heeft gelijke zijden, gelijke hoeken en heeft vijf diagonalen.}
    \antwoord{Deze uitspraak bevat overbodige informatie.}

\end{vraagenantwoord}

\afbeelding{img/veelhoeken.jpg}{}{}

\vraag{Formuleer de definitie van een trapezium, een parallellogram, een regelmatige veelhoek en een cirkel.
Onderzoek telkens welke van de onderstaande figuren voldoen aan deze definitie en welke niet.}
\end{vraagenantwoord}  

\afbeelding[13.2cm]{img/loepEls/50vlakkies.png}{}{}

\begin{vraagenantwoord}[resume]
    

\vraag{\textbf{Extra:} Onderzoek of we de onderstaande uitspraak zouden kunnen gebruiken als definitie voor een vierkant.
Zo nee, kun je hier informatie aan toevoegen/weglaten zodat het wel kan?
\begin{citaat}
    Een vierkant is een vierhoek waarvan de diagonalen loodrecht staan en even lang zijn.
\end{citaat}}
\antwoord{Een vierkant is een vierhoek waarvan de diagonalen loodrecht staan, elkaar middendoor delen en even lang zijn.}

\end{vraagenantwoord}
\end{lesactiviteit}

\begin{multicols}{2}
De extra opgave is een mooie uitdaging om in de leerstof vooruit te kijken.
Ze geeft leerlingen, die daaraan toe zijn, een voorproefje van de volgende niveaus van de theorie van Van Hiele.

\subsection{Redeneringen staven met informele argumenten}\label{redeneringen-staven-met-informele-argumenten}

\subsubsection*{Mind the gap}

Wanneer we redeneringen met meetkundige figuren opbouwen, gebruiken we vaak een prototype als denkmodel.
Een prototype is een voorbeeld van de figuur in zijn meest willekeurige vorm.
Zo gebruiken we voor een redenering met driehoeken, geen gelijkzijdige driehoek maar werken we met een willekeurige driehoek zonder speciale eigenschappen.
We willen immers een redenering die algemeen geldig is.

Leerlingen zijn zich vaak niet bewust van dat uitgangspunt.
Bovendien werken opgaven in handboeken dit soms tegen.
In \autoref{fig:opperdrie} zie je hoe de oppervlakteformule voor een willekeurig driehoek met behulp van een rechthoekige driehoek wordt afgeleid door de driehoek om te structureren naar een rechthoek.
Deze redenering verschilt voor andere, niet rechthoekige driehoeken.
Daar structureer je op een andere manier om naar een rechthoek (zie \autoref{fig:driehoekomvormen}) of, zelfs eenvoudiger naar een parallellogram.
\end{multicols}
\afbeelding[12cm]{img/loepEls/image6.png}{Oppervlakte driehoek (Uit: \emph{Wiskidz 5}, blok 9, p. 12)}{fig:opperdrie}
\begin{multicols}{2}

\afbeelding[3cm]{img/loepEls/driehoekomvormen.png}{Driehoek reconstrueren tot rechthoek}{fig:driehoekomvormen}


\subsubsection*{Hoe bouw je de brug?}

In de eerste plaats is het belangrijk dat leerlingen begrijpen dat je op basis van één voorbeeld een redenering niet kunt veralgemenen.
In principe ben je er maar zeker van dat een redenering klopt, als je ze gecontroleerd hebt voor \textbf{alle} elementen van de verzameling (of als je ze bewezen hebt, cfr. niveau 4).
We kunnen echter niet alle elementen effectief controleren.
Daarom gebruiken we op niveau 3 meer informele strategieën.
Bijvoorbeeld: we checken de uitspraak voor verschillende types voorbeelden of redeneren op een algemeen prototype.
Deze manier van werken helpt leerlingen om te groeien in het concept van bewijzen.
De volgende lesactiviteit illustreert hoe je dit voor de formule van de oppervlakte van een driehoek kunt aanpakken.
\end{multicols}
\begin{lesactiviteit}{De oppervlakte van een driehoek}
Bij het bestuderen van driehoeken, kunnen we ze indelen volgens de hoeken (stomp-, recht- en scherphoekig) of volgens de zijden (gelijkzijdig, gelijkbenig en ongelijkbenig).
Wanneer je beide criteria combineert, levert dat maximaal 9 verschillende combinaties op.

\begin{vraagenantwoord}
    \vraag{Teken voor elke combinatie een driehoek.\label{q:alledriehoeken}}
    \antwoord{Dit levert slechts 7 verschillende driehoeken op (zie de onderstaande figuur).
    Een gelijkzijdige driehoek is altijd scherphoekig.
    \afbeelding{img/loepEls/alledriehoeken.png}{}{}
    }

    \vraag{We zoeken een formule voor de oppervlakte van een driehoek.
    Structureer daarvoor de onderstaande driehoek om tot een figuur waarvoor je al een oppervlakteformule kent.
    Bepaal zo de formule voor een driehoek.
    \afbeelding[3cm]{img/loepEls/willekeurigedriehoek.png}{}{}
    }
    \antwoord{
    \afbeelding[12cm]{img/loepEls/driehoektopara.png}{}{}
    In de bovenstaande figuur definiëren we de hoogte van de driehoek als de afstand van de top tot de overstaande zijde, die we de basis noemen.
    Vervolgens leid je de oppervlakteformule af:
    \begin{align*}
        \text{oppervlakte driehoek} &= \frac{\text{oppervlakte parallellogram}}{2}\\
        &= \frac{\text{basis}_{\text{parallellogram}}\text{ × }\text{hoogte}_{\text{parallellogram}}}{2}\\
        &= \frac{\text{basis}_{\text{driehoek}}\ \text{× hoogte}_{\text{driehoek}}}{2}\\
    \end{align*}}

    \vraag{We willen deze formule veralgemenen voor alle driehoeken.
    Controleer of deze redenering geldig is voor de andere driehoeken die je tekende in \autoref{q:alledriehoeken}.}
    \antwoord{Je kunt elke driehoek omstructureren naar een parallellogram.
    Voor de rechthoekige driehoeken levert dit een rechthoek op, wat ook een parallellogram is.
    Zo veralgemeen je de redenering.}

    \vraag{Als je een andere zijde als basis neemt, kun je dan nog steeds deze redenering gebruiken?
    Controleer voor enkele driehoeken.}
    \antwoord{Je kunt de driehoek steeds omvormen naar een parallellogram, met bijbehorende basis en hoogte.
    Het is zinvol om de verschillende hoogtes en basissen te laten aanduiden.
    Leerlingen hebben het soms moeilijk om hoogtes loodrecht aan te duiden.
    Rechthoekige en stomphoekige driehoeken verdienen hierbij zeker bijzondere aandacht omdat basis en hoogte dan respectievelijk samenvallen met de zijden of de hoogte buiten de driehoek ligt.}
\end{vraagenantwoord}
\end{lesactiviteit}
\begin{multicols}{2}
Op basis van dit onderzoek besluiten we dat de gevonden formule voor de oppervlakte van een driehoek algemeen geldig is.

\subsection{Classificatie van meetkundige objecten}\label{classificatie-van-meetkundige-objecten}

\subsubsection*{Mind the gap}
Sinds de invoering van de eindtermen in 1997, kregen de verzamelingenleer en de logica in het basisonderwijs een ander statuut:

\begin{citaat}
Het inventief en inzichtelijk werk van kinderen kan niet starten vanuit een opgelegd abstract raamwerk, toch niet in eerste instantie.
Vandaar dat het begrippenarsenaal uit de verzamelingenleer niet meer als doel op zich in de eindtermen voorkomt, al kunnen sommige voorstellingswijzen (venndiagrammen, relatiepijlen, ...) interessante hulpmiddelen blijven voor het wiskundig denken van kinderen. (AHOVoKS, sd))
\end{citaat}

Deze, overigens waardevolle, nuancering heeft er toe geleid dat het gebruik van verzamelingen als interessante voorstellingswijze vrijwel geheel uit het basisonderwijs is verbannen.
Handboekmakers zoeken krampachtig naar andere middelen om verbanden tussen de verschillende meetkundige objecten duidelijk te maken.
Een voorbeeld van zo'n schema vind je in \autoref{fig:catego}.

Deze figuur is niet alleen een zeer complex boomdiagram maar bevat foute informatie die leerlingen op het verkeerde been zet.
Een boomstructuur is nuttig en handig als de verschillende takken van de boom gescheiden categorieën bevatten.
Voor het classificeren van vlakke figuren is dit niet zo: bv. een vierkant is ook een rechthoek, ruit
\ldots{}
In het bovenstaande schema staat het vierkant in een andere tak dan de rechthoek en lijkt er zelfs helemaal geen verband met de tak van de ruiten.
Het gebruik van venndiagrammen zou in dit geval de verbanden tussen de figuren beter illustreren.

In het secundair onderwijs maakt de verzamelingenleer wel deel uit van het programma.
Dit onderdeel wordt vaak als een apart hoofdstuk behandeld.
Zoals je in \autoref{fig:vennie} ziet, leidt dit soms tot wat artificiële contexten om bijvoorbeeld het begrip `deelverzameling' betekenis te geven.

\end{multicols}

\afbeelding[14cm]{img/loepEls/image7.png}{Meetkundig boomdiagram (Uit: \emph{Wiskidz 6}, blok 4, les 5)}{fig:catego}

\begin{multicols}{2}
    


\end{multicols}

\afbeelding[14cm]{img/loepEls/image8.png}{Verzamelingenleer in de meetkunde (Uit: \emph{Delta Nova 1}, H1, p. 24)}{fig:vennie}

\begin{multicols}{2}
    

\subsubsection*{Hoe bouw je de brug?}

Het (functioneel) gebruiken van verzamelingen/venndiagrammen bij meetkundige inzichten kan op dit vlak een win-win situatie opleveren.
\autoref{fig:vennexample} toont hoe verzamelingen enerzijds een krachtig hulpmiddel zijn om voor meetkundige begrippen tot een goede begripsinvulling te komen.
Anderzijds is meetkunde concreet genoeg om het abstractere denken van de verzamelingenleer te concretiseren.

\afbeelding[14cm]{img/loepEls/vennexample.png}{Classificatie rechthoeken}{fig:vennexample}

Dat zou misschien tegemoet komen aan de verwarring van de leerling die we zien in \autoref{fig:rechtclass}.
Ze worstelt nog met de betekenis van het begrip rechthoek.
Daarom schrijft ze naast de `rechthoeken', eerst het begrip `vierkant' en streept dit later opnieuw door wanneer ze beseft dat een vierkant ook een rechthoek is.



\subsection{Eigenschappen en eenvoudige als-dan-redeneringen}\label{eigenschappen-en-eenvoudige-als-dan-redeneringen}

\subsubsection*{Mind the gap}

Een eigenschap van een meetkundige figuur is een uitspraak die voor al de vertegenwoordigers van die figuur klopt.
Qua redenering loopt dit analoog met wat we al voor definities formuleerden.
In het basisonderwijs onderzoeken leerlingen eigenschappen op basis van voorbeelden.
Een typische opgave vind je in \autoref{fig:eigenschappen}.


In deze opgave onderzoeken de leerlingen verschillende elementen (hoeken, lengtes van zijden\ldots) in het getekende vierkant en de rechthoek.
De opdracht is vooral uitvoerend en blijft beperkt tot de constatatie van wat er voor deze (voorbeeld)figuur klopt.
Er is geen opstapje naar een veralgemening, een uitgebreidere formulering van de eigenschappen of de meerwaarde van een eigenschap.


\end{multicols}

\afbeelding[14cm]{img/loepEls/image9.png}{Andere classificatie rechthoeken (Uit: \emph{Nando 1}, H13, p. 10)}{fig:rechtclass}

\afbeelding[14cm]{img/loepEls/image10.png}{Eigenschappen van rechthoek en vierkant (Uit: \emph{Wiskanjers 4}, blok 3, les 5)}{fig:eigenschappen}

\begin{multicols}{2}
    

\subsubsection*{Hoe bouw je de brug?}

Naast het controleren van eigenschappen op verschillende types van figuren, zoals we in de paragraaf over \textbf{betekenisvolle definities} al toelichtten, kan het \textbf{zelf} tekenen van figuren helpen om de meerwaarde van een eigenschap te ontdekken.

Zo is bv. de symmetrieas een handig hulpmiddeltje om de tophoek van een gelijkbenige driehoek goed te plaatsen.
Bij een gelijkzijdige driehoek gebruik je die, in combinatie met een passer, om zo de tophoek te vinden.
Of bij het tekenen van een ruit is het handig om te vertrekken van de diagonalen die loodrecht moeten zijn en elkaar middendoor moeten snijden\ldots{}
Al deze hulpmiddelen steunen op eigenschappen van deze figuren.

Constructieactiviteiten, die de laatste jaren minder aandacht krijgen, zijn rijke opdrachten die ervoor zorgen dat meetkundige figuren met hun eigenschappen beter worden verkend.
Het opent later de deur naar het maken van een snelle schets of een tekening wanneer je een uitspraak moet beoordelen.
Een voorbeeld van zo'n opgave vind je in de lesactiviteit hieronder.
\end{multicols}

\begin{lesactiviteit}{Construeren en als-dan-uitspraken beoordelen}
\begin{vraagenantwoord}
    \vraag{Lees even onderstaande uitspraken, en bedenk of je ermee akkoord gaat of niet.
    \begin{itemize}
    \item Als in een vierhoek de diagonalen elkaar middendoor snijden en loodrecht zijn, dan is de vierhoek een ruit.
    \item Als de diagonalen van een vierhoek even lang zijn, dan is de vierhoek een rechthoek.
    \item Als de diagonalen van een parallellogram even lang zijn, dan is dit parallellogram een vierkant.
    \end{itemize}
    }
    \antwoord{Het is niet altijd gemakkelijk om te beoordelen of een uitspraak klopt.
    Wanneer je één voorbeeldfiguur kunt tekenen waarvoor de uitspraak niet klopt, dan spreken we van een onware uitspraak.
    Om te controleren dat de uitspraak waar is, zou je in principe alle figuren moeten onderzoeken.
    Dat kan natuurlijk niet.
    Daarom bouwen we een redenering op waarmee we verklaren waarom een uitspraak waar is.
    In dat geval noemen we deze uitspraak `een eigenschap'.
    Het tekenen van verschillende voorbeeldfiguren kan helpen om dit soort redeneringen op te bouwen.
    }

    \vraag{Teken in elk vakje in de onderstaande tabel een figuur als dat kan.
    Maak gebruik van je passer en geodriehoek.    
    In een bepaalde kolom (bv. de kolom ``Trapezium'') teken je enkel die figuren die niet aan een strenger criterium voldoen (dus geen parallellogram, ruit, rechthoek of vierkant, ook al zijn dit allemaal trapezia).
    De symbolen in de kolom ``Diag.'' betekenen hetvolgende:
    \begin{align*}
        = &: \text{diagonalen zijn even lang,}\\
        \perp &: \text{diagonalen staan loodrecht op elkaar,}\\
        X &: \text{diagonalen snijden elkaar middendoor.}
    \end{align*}
    Indien er een `\(\text{niet}\)' voor het symbool staat, betekent dit dat de tegenstelde uitspraak bedoeld wordt. 
    Bijvoorbeeld `\(\text{niet }\perp\)' betekent dat de diagonalen niet loodrecht staan.
    }

    \begin{tabel}
	\begin{tabular}{|>{\raggedright}m{1cm}|m{2cm}|m{1.6cm}|m{2.3cm}|m{1cm}|m{1.6cm}|m{1.5cm}|}
		\hline
		\tableheader{Diag.} & \tableheader{Willekeurige vierhoek} & \tableheader{Trapezium} & \tableheader{Parallellogram} & \tableheader{Ruit} & \tableheader{Rechthoek} & \tableheader{Vierkant} \\ 
		\hline
        \(=\)\newline\(\perp\)\newline\(X\) & & & & & & \\
        \hline
        \(=\)\newline\(\perp\)\newline\(\text{niet }X\) & & & & & & \\
        \hline
        \(=\)\newline\(\text{niet }\perp\)\newline\(X\) & & & & & & \\
        \hline
        \(=\)\newline\(\text{niet }\perp\)\newline\(\text{niet }X\) & & & & & & \\
        \hline
        \(\text{niet }=\)\newline\(\perp\)\newline\(X\) & & & & & & \\
        \hline
        \(\text{niet }=\)\newline\(\perp\)\newline\(\text{niet }X\) & & & & & & \\
        \hline
        \(\text{niet }=\)\newline\(\text{niet }\perp\)\newline\(X\) & & & & & & \\
        \hline
        \(\text{niet }=\)\newline\(\text{niet }\perp\)\newline\(\text{niet }X\) & & & & & & \\
        \hline
	\end{tabular}
	\caption{Welke vierhoeken zijn mogelijk met (on)gelijke, (niet-)loodrechte, (niet-)gelijkgedeelde diagonalen?}
    \label{tab:eigenschapmeetkunde}
    \end{tabel}

    \antwoord{Je kan een figuur tekenen in de kolom van
    \begin{itemize}
    \item Rij 1: een vierkant
    \item Rij 2: een willekeurige vierhoek, een (gelijkbenig) trapezium
    \item Rij 3: een rechthoek
    \item Rij 4: een willekeurige vierhoek, een (gelijkbenig) trapezium
    \item Rij 5: een ruit
    \item Rij 6: een willekeurige vierhoek (vlieger), een trapezium
    \item Rij 7: een parallellogram
    \item Rij 8: een willekeurige vierhoek, een trapezium)
    \end{itemize}
    }

    \vraag{Beoordeel nu volgende uitspraken.
    Gebruik hiervoor de tekeningen uit de vorige tabel.}
    \begin{vraagenantwoord}
        \vraag{Als in een vierhoek de diagonalen elkaar middendoor snijden en loodrecht zijn, dan is deze vierhoek een ruit.}
        \antwoord{Waar.}
        \vraag{Als de diagonalen van een vierhoek even lang zijn, dan is de vierhoek een rechthoek.}
        \antwoord{Niet waar, het kan ook een trapezium of willekeurige vierhoek zijn.}
        \vraag{Als de diagonalen van een parallellogram even lang zijn, dan is dit parallellogram een vierkant.}
        \antwoord{Niet waar, ook een rechthoek voldoet hieraan.}
    \end{vraagenantwoord}

\end{vraagenantwoord}   
\end{lesactiviteit}

\begin{multicols}{2}
\section{Meten: herleidingen}\label{maten}

\subsection{Leerstof in het basisonderwijs}



Het herleiden van metingen in de ene maat, naar een resultaat in een andere maat gebeurt in het basisonderwijs in de meeste gevallen via een tabel zoals in \autoref{fig:omzettingstabel}.
In het beste geval leren leerlingen niet alleen \textbf{hoe} zo'n tabel werkt maar ook \textbf{waarom} de werkwijze werkt.
Hieronder maken we dat inzicht helder door de achterliggende redenering te expliciteren.
We redeneren voor het volgende voorbeeld:
\[
\SI{567}{\deci\metre\squared} = \ldots\ldots\ldots\si{\metre\squared}
\]
\end{multicols}
\afbeelding[14cm]{img/loepEls/image11.png}{Omzettingstabel voor oppervlaktematen (Uit: \emph{Wiskidz 6}, blok 7, p. 9)}{fig:omzettingstabel}

\begin{multicols}{2}
    

We steunen hierbij op twee basisinzichten voor meten:
\begin{itemize}
\item De verhouding tussen opeenvolgende maten is constant. Voor oppervlakte past een kleinere maat \(100\) keer in een grotere maat.
\item Maat en maatgetal zijn omgekeerd evenredig.
Als we een oppervlakte uitdrukken in een kleinere maat wordt het maatgetal groter.
\end{itemize}

\begin{enumerate}
\item We bouwen de tabel voor oppervlaktematen op:
\begin{tabel}
    \begin{tabular}{|m{0.3cm}:m{0.3cm}|m{0.3cm}:m{0.3cm}|m{0.3cm}:m{0.3cm}|m{0.3cm}:m{0.3cm}|m{0.3cm}:m{0.3cm}|m{0.3cm}:m{0.3cm}|}
		\hline
        \multicolumn{2}{|c}{\(\si{\kilo\metre\squared}\)} & \multicolumn{2}{|c}{\(\si{\hecto\metre\squared}\)} & \multicolumn{2}{|c}{\(\si{\deca\metre\squared}\)} & \multicolumn{2}{|c}{\(\si{\metre\squared}\)} & \multicolumn{2}{|c}{\(\si{\deci\metre\squared}\)} & \multicolumn{2}{|c|}{\(\si{\centi\metre\squared}\)} \\
        \hline
	\end{tabular}
\end{tabel}

\item \emph{De verhouding tussen twee opeenvolgende maten bij oppervlakte is \(100\)} (basisinzicht).
Dat betekent dat er \(\SI{100}{\deci\metre\squared}\) passen in \(\SI{1}{\metre\squared}\).
We moeten m.a.w. \(99\) in de kolom van de vierkante decimeter kunnen schrijven voor we cijfers in de kolom van de vierkante meter invullen.
Dat betekent dat er bij elke maat plaats voorzien wordt voor twee cijfers.
Daarom verdelen we elke kolom onder in twee subkolommen.

\begin{tabel}
	\begin{tabular}{|m{0.3cm}:m{0.3cm}|m{0.3cm}:m{0.3cm}|m{0.3cm}:m{0.3cm}|m{0.3cm}:m{0.3cm}|m{0.3cm}:m{0.3cm}|m{0.3cm}:m{0.3cm}|}
		\hline
        \multicolumn{2}{|c}{\(\si{\kilo\metre\squared}\)} & \multicolumn{2}{|c}{\(\si{\hecto\metre\squared}\)} & \multicolumn{2}{|c}{\(\si{\deca\metre\squared}\)} & \multicolumn{2}{|c}{\(\si{\metre\squared}\)} & \multicolumn{2}{|c}{\(\si{\deci\metre\squared}\)} & \multicolumn{2}{|c|}{\(\si{\centi\metre\squared}\)} \\
        \hline
        .&.&.&.&.&.&.&.&.&.&.&.\\
        \hline
	\end{tabular}
\end{tabel}

\item We noteren de meting, \(\SI{567}{\deci\metre\squared}\), in de tabel.
We weten dat de maat (\(\si{\deci\metre\squared}\)) slaat op het cijfer van de eenheden.
De \(7\) hoort dus op de plaats van de eenheden bij de \(\si{\deci\metre\squared}\).

\begin{tabel}
    \begin{tabular}{|m{0.3cm}:m{0.3cm}|m{0.3cm}:m{0.3cm}|m{0.3cm}:m{0.3cm}|m{0.3cm}:m{0.3cm}|m{0.3cm}:m{0.3cm}|m{0.3cm}:m{0.3cm}|}
		\hline
        \multicolumn{2}{|c}{\(\si{\kilo\metre\squared}\)} & \multicolumn{2}{|c}{\(\si{\hecto\metre\squared}\)} & \multicolumn{2}{|c}{\(\si{\deca\metre\squared}\)} & \multicolumn{2}{|c}{\(\si{\metre\squared}\)} & \multicolumn{2}{|c}{\(\si{\deci\metre\squared}\)} & \multicolumn{2}{|c|}{\(\si{\centi\metre\squared}\)} \\
        \hline
        .&.&.&.&.&.&.&5&6&7&.&.\\
        \hline
	\end{tabular}
\end{tabel}

\item We willen het oorspronkelijke meetresultaat omzetten naar \(\si{\metre\squared}\).
\emph{De nieuwe maat is groter dan de oorspronkelijke maat.
Het maatgetal zal dus kleiner zijn} (basisinzicht).

\item Het nieuwe cijfer van de eenheden wordt dat wat er bij de eenheden van de vierkante meter staat.
Dit is de \(5\).
We plaatsen een komma achter deze \(5\), want de maat moet op het cijfer van de eenheden slaan.

\item Antwoord: \(\SI{567}{\deci\metre\squared} = \SI{5.67}{\metre\squared}\)
\end{enumerate}

Door het benadrukken van de basisinzichten, maak je het mogelijk om, na verloop van tijd, over te schakelen op een pijlenschema en later op een abstracte redenering (zoals onderaan in \autoref{fig:omzettingstabel}).

\end{multicols}

\afbeelding[14cm]{img/loepEls/image12.png}{Herleidingsoefeningen (Uit: \emph{Delta Nova 1}, p. 388)}{fig:herleiding}

\begin{multicols}{2}
    
\subsection{Leerstof in het secundair onderwijs}



In het secundair onderwijs streven we ernaar dat leerlingen herleidingen zoals in \autoref{fig:herleiding} abstract leren uitvoeren.
Dit steunt op dezelfde basisinzichten.
We geven een voorbeeld hoe zo'n abstracte omzetting inzichtelijk verloopt:
\[
\SI{45375}{\centi\metre\squared} = \ldots\ldots\ldots\si{\deci\metre\squared}
\]
\begin{itemize}
    \item Welke eenheid is het grootst? De \(\si{\deci\metre\squared}\).
    We herleiden van een kleinere maat naar een grotere maat, het maatgetal zal dus kleiner worden (basisinzicht).
    \item De verhouding tussen opeenvolgende lengtematen is \(100\) (basisinzicht).
    De maat wordt \(100\) keer groter, het maatgetal \(100\) keer kleiner.
    In dit geval betekent dat dus:
    \[
    \SI{1}{\centi\metre\squared} = \frac{1}{100}\si{\deci\metre\squared} = \SI{0.01}{\deci\metre\squared}
    \]
    Toegepast op het voorbeeld, krijgen we:
    \begin{align*}
        \SI{45375}{\centi\metre\squared} &= 45~375\text{ keer } \left(\SI{1}{\centi\metre\squared}\right)\\
        &= 45~375\text{ keer } \left(\SI{0.01}{\deci\metre\squared}\right)\\
        &= 45~375 \cdot\SI{0.01}{\deci\metre\squared}\\
        &= \SI{453.75}{\deci\metre\squared}
    \end{align*}
\end{itemize}
\subsection{Mind the gap}

\subsubsection*{Onvoldoende vertrouwdheid met het metriek stelsel}\label{onvoldoende-vertrouwdheid-met-het-metriek-stelsel}

In het basisonderwijs gebruiken we de termen maat en maatgetal, terwijl we in het secundair onderwijs vaker de term (maat)eenheid en eenheden hanteren.

De voorvoegsels `hecto' en `deca' worden maar zelden aangeleerd.
Zoals je op \autoref{fig:omzettingstabel} ziet, noteren ze de hectometer en de decameter met \(\SI{100}{\metre}\) en \(\SI{10}{\metre}\).
Daardoor verliest het metriek stelsel aan systematiek.
Het verband met de tientalligheid van ons talstelsel en de systematiek van het metriek stelsel komt veel minder duidelijk tot uiting.

Het betekent dat leerlingen met deze maateenheden ook weinig of geen concrete ervaringen hebben opgedaan.
Daardoor kunnen ze minder goed inschatten hoe groot deze maten zijn en of hun berekeningen realistisch zijn.
Hopelijk komt hierin verandering met de nieuwe minimumdoelen voor het basisonderwijs.

\subsubsection*{Het veralgoritmiseren van herleiden met behulp van de tabel}\label{het-veralgoritmiseren-van-herleiden-met-behulp-van-de-tabel}

In handboeken van zowel het basisonderwijs (in \autoref{fig:herlbasis}) als van het secundair onderwijs (in \autoref{fig:herlsecund}) zien we vaak hoe het omzetten tot \textbf{een techniek} wordt herleid.
Een stappenplan of `algoritme' werkt het gedachteloos uitvoeren van deze techniek in de hand.
Met allerlei trucjes worden nullen bijgeplaatst, komma's verschoven, kolommen ingevoerd.
Het aanleren van deze `regeltjes' legt krediet op het doorgroeien naar het abstract uitvoeren van herleidingen.
Je maakt kinderen m.a.w. afhankelijk van het algoritme, in plaats van een basis te geven om later, vb. bij fysica, sneller abstract herleidingen uit te voeren.
\end{multicols}

\afbeelding[14cm]{img/loepEls/image13.png}{Algoritme voor eenheden herleiding basisonderwijs (Uit: \emph{Wiskidz 6}, blok 7, p. 9)}{fig:herlbasis}

\afbeelding[14cm]{img/loepEls/image14.png}{Algoritme voor eenheden herleiding secundair onderwijs (Uit: \emph{Nando 1}, H13, p. 15)}{fig:herlsecund}

\begin{multicols}{2}

Niet alle handboeken bouwen deze redeneringen op dezelfde manier op.
In \autoref{fig:herlsecund2} pakt men het op een heel andere manier aan.

\end{multicols}

\afbeelding[14cm]{img/loepEls/image15.png}{Een andere methode uit het secundair onderwijs (Uit: \emph{Delta Nova 1}, H13, p. 387)}{fig:herlsecund2}

\begin{multicols}{2}

In deze redenering wordt met pijlenschema's gewerkt.
Dit soort schema's ondersteunt het begrijpen van recht- en omgekeerd evenredige grootheden.
Er wordt hier echter stilgestaan bij het rekenwerk (de vermenigvuldiging van \(\SI{1}{\kilo\metre}\) naar \(\SI{0.002516}{\kilo\metre}\) en niet bij het inzicht (de omgekeerd evenredigheid van maat en maatgetal).

\subsection{Hoe bouw je de brug}

In een tussenfase kun je een pijlenschema gebruiken om de brug te maken.
In sommige basisscholen zullen leerlingen dit al geleerd hebben (zoals in \autoref{fig:omzettingstabel}), maar vaak nog niet.
We leggen de opbouw uit aan de hand van een voorbeeld voor oppervlakte.
\[
\SI{1325}{\deci\metre\squared} = \ldots\ldots\ldots\si{\metre\squared}
\]

\begin{itemize}
\item Deze herleiding gaat over oppervlaktematen.
De verschillende oppervlaktematen zijn:
\begin{tabel}
    \begin{tabular}{|m{0.7cm}|m{0.7cm}|m{0.7cm}|m{0.7cm}|m{0.7cm}|m{0.7cm}|m{0.7cm}|}
		\hline
        \(\si{\kilo\metre\squared}\) & \(\si{\hecto\metre\squared}\) & \(\si{\deca\metre\squared}\) & \;\(\si{\metre\squared}\) & \(\si{\deci\metre\squared}\) & \(\si{\centi\metre\squared}\) & \(\si{\milli\metre\squared}\)\\
        \hline
	\end{tabular}
\end{tabel}
\(\SI{10000}{\metre\squared}\) noemen we een vierkante hectometer, \(\SI{100}{\metre\squared}\) een vierkante decameter.

\item We weten dat bij de oppervlaktematen \emph{de verhouding tussen opeenvolgende maten een factor \(100\) is} (basisinzicht).
Dat betekent dat \(\SI{1}{\metre\squared} = \SI{100}{\deci\metre\squared}\).
Of dat er in één vierkante meter honderd vierkante decimeters passen.
In \autoref{fig:verband} zie je de verbanden tussen \emph{de opeenvolgende maten}.

\item De maat(eenheid) en het maatgetal (of het aantal eenheden) zijn omgekeerd evenredig.
Dat betekent: hoe groter de maat waarmee je meet, hoe kleiner het maatgetal en omgekeerd.
Een oppervlakte in vierkante meter uitdrukken, zal een kleiner getal opleveren dan als je in vierkante decimeter meet.

\item Een pijlenschema helpt om de omzetting te maken.
Dit pijlenschema (in \autoref{fig:herleidenbeter}) verschilt fundamenteel van het schema uit \autoref{fig:herlsecund2}, omdat hier de nadruk op de omgekeerde evenredigheid van maat en maatgetal en de verhouding tussen de opeenvolgende maateenheden wordt gelegd.

\end{itemize}

\end{multicols}

\afbeelding[12cm]{img/loepEls/verbandenoppervlaktematen.jpg}{Verbanden tussen oppervlaktematen}{fig:verband}

\afbeelding[14cm]{img/loepEls/pijlenomzetting.png}{Alternatief pijlenschema voor herleidingen}{fig:herleidenbeter}



\begin{multicols}{2}
    

Naarmate leerlingen vlotter worden in het uitvoeren van de bovenstaande werkwijze, kun je die verkorten tot de onderstaande notatie.

\afbeelding[6cm]{img/loepEls/verkortherleiden.png}{Verkort pijlenschema voor herleiden}{fig:verkortherleiden}

In een latere fase laten we de pijlen achterwege en redeneren we volledig abstract.

\end{multicols}

%% BRONNEN %%%%%%%%%%%%%%%%%%%%%%%%%%%%%%%%%%%%%%%%%%%%%%%%%%%%%%%%%%%%%%%%%%%%%
\begin{bronnen}
    \item AHOVoKS, A. v. (Red.). (z.d.). \emph{Onderwijsdoelen - Uitgangspunten}. Opgeroepen op maart 10, 2024, van \url{https://www.onderwijsdoelen.be/uitgangspunten/4680}
    \item Carmeliet, C. D., Serneels, K., \& Tytgat, P. (2020). \emph{Delta nova 1a \& b.} Antwerpen: Plantyn.
    \item Carmeliet, C., Deloddere, N., Serneels, K., \& Tytgat, P. (z.d.). \emph{Delta Nova 1.} Plantyn. Opgeroepen op maart 2024, van \url{https://view.publitas.com/plantyn/delta-nova-1-digitale-licentie/page/1}
    \item Carmliet, C., Deloddere, N., Serneels, K., \& Tytgat, P. (z.d.). \emph{Delta Nova 1.} Plantyn. Opgeroepen op maart 2024, van \url{https://view.publitas.com/plantyn/delta-nova-1-digitale-licentie/page/1}
    \item Carreyn, B., Geeurickx, F., \& Van Nieuwenhuyze, R. (2022). \emph{Nando 1.} Brugge: Die Keure.
    \item Carreyn, B., Geeurickx, F., \& Van Nieuwenzhuyze, R. (2022). \emph{Nando 2.} Brugge: Die Keure.
    \item Carreyn, B.; Geeurickx, F., \& Van Nieuwenhuyze, R. (z.d.). \emph{Van basis tot limiet 1}, Brugge: Die Keure.
    \item Denys, D., Vanpaemel, A., \& e.a.. (2022). \emph{Reken Maar.} Wommelgem: Van In.
    \item De Wilde, C., Libbrecht, S., \& Mivis, K. (z.d.). \emph{Wiskidz} (Vol. zesde leerjaar blok 2 (p 15)).
    \item De Wilde, C., Libbrecht, S., Mivis, K., \& e.a.. (z.d.). \emph{Wiskidz.} Averbode: Uitgeverij Averbode.
    \item Duerloo, M., Gobien, S., \& e.a.. (2022). \emph{Wiskanjers.} Antwerpen: Plantyn.
    \item Thaels, K., Eggermont, H., \& Janssens, D. (2001). \emph{Van ruimtelijk inzicht naar ruimtemeektunde.} Deurne: Wolters Plantyn.
\end{bronnen}